% ============================================================
%  ch4_methodology.tex
%  Chapter 4: Methodology
%  Master's Thesis -- SQL MATCH_RECOGNIZE on Pandas DataFrames
% ============================================================
% Provide \mr if not already defined by the main document
\providecommand{\mr}{\texttt{MATCH\_RECOGNIZE}}
\chapter{Methodology}
\label{chap:methodology}

Our \texttt{MATCH\_RECOGNIZE} implementation follows a modular,
five-stage architecture specifically designed for SQL row-pattern
recognition.  The engine executes SQL queries directly on
\texttt{pandas} DataFrames~\cite{monierashraf2025rowmatch} while
fully preserving SQL:2016 semantics~\cite{iso2016,ISO2021,melton2017sqlrpr}.
The pipeline consists of five phases: \emph{parsing},
\emph{semantic validation}, \emph{logical plan construction},
\emph{pattern compilation}, and \emph{execution}
(Fig.~\ref{fig:system-architecture}).

In the parsing stage, the SQL query is transformed into an abstract
syntax tree (AST) that captures both the \texttt{SELECT}/\texttt{FROM}
clauses and \texttt{MATCH\_RECOGNIZE} sub-clauses~\cite{tarazona2021antlr,parr2013definitive}.
Semantic validation then checks schema references, variable usage,
predicates, skip policies, and output modes.  Compliance rules are
applied to ensure standard conformance, correct use of navigation
functions, consistent measure semantics, and robust handling of edge
cases such as empty matches, partition boundaries, and skip
behaviours.

The validated plan is translated into a logical plan composed of
modular operators for partitioning, sorting, matching, measuring,
and skipping.  From this plan, the system compiles a
nondeterministic finite automaton (NFA) with guarded transitions
that support quantifiers, alternation, and \texttt{PERMUTE}
constructs~\cite{salomaa1996nfa}.  The NFA is constructed using
Thompson's construction, extended for SQL:2016 semantics by
embedding predicates from the \texttt{DEFINE} clause, minimising
$\varepsilon$-transitions, and applying subset construction with
state deduplication to mitigate exponential growth.

Finally, execution is performed over ordered partitions using the
compiled automaton.  Matches are materialised with the correct
\texttt{RUNNING}/\texttt{FINAL} semantics.  The evaluation engine
manages context separation, applies navigation functions
(\texttt{FIRST}, \texttt{LAST}, \texttt{PREV}, \texttt{NEXT}) with
partition awareness, and computes aggregates such as \texttt{COUNT},
\texttt{SUM}, and \texttt{AVG} across pattern variables.  This
ensures full compliance with SQL:2016 while producing results
directly as \texttt{pandas} DataFrames.

The remainder of this chapter is organised as follows.
Section~\ref{sec:meth-overview} provides a high-level overview of
the three primary runtime components.
Section~\ref{sec:meth-compiler} describes the SQL parsing pipeline
in detail.
Section~\ref{sec:meth-semantic} covers the semantic validation
checks applied after parsing.
Section~\ref{sec:meth-logical} presents the logical plan
construction step.
Section~\ref{sec:meth-automata} establishes the finite-automata
foundations.
Section~\ref{sec:meth-pandas} presents the Pandas execution model.
Section~\ref{sec:meth-impl} summarises the concrete implementation.
Section~\ref{sec:bg-performance} details the
performance-optimisation layer.

% ----------------------------------------------------------------
\section{System Architecture Overview}
\label{sec:meth-overview}
% ----------------------------------------------------------------

The engine's five-phase pipeline is designed so that each stage
produces a well-defined artefact consumed by the next.
The following three subsections give a brief description of the
primary runtime components before the detailed treatments in
Sections~\ref{sec:meth-compiler}--\ref{sec:meth-pandas}.

% ................................
\subsection{SQL Parsing and AST Construction}
\label{subsec:overview-parsing}
% ................................

The process starts by analysing the user's SQL query and identifying
its logical components.  A parsing process explores the query,
recognises its main components---grouping (\texttt{PARTITION BY}),
ordering (\texttt{ORDER BY}), the pattern to be searched, and any
specific conditions or output metrics---and organises these into an
Abstract Syntax Tree (AST).  The AST provides a clear, structured
representation of the query.  This approach eases query validation
and prepares it for pattern matching in the following stages.

% ................................
\subsection{Pattern Translation to Finite Automata}
\label{subsec:overview-automata}
% ................................

Once the AST is complete, the system focuses on the pattern
component of the SQL query.  The pattern is converted into a finite
automaton that defines a set of states and transitions for detecting
complex sequences or behaviours in the data.  The process begins by
constructing a non-deterministic finite automaton (NFA) from the
pattern and then, for improved efficiency, transforms it into a
deterministic finite automaton
(DFA)~\cite{georgiou2025streaming}.  The DFA eliminates transition
ambiguity, enabling the system to scan data sequences quickly and
accurately.  This translation handles a wide variety of pattern
features---repetitions, alternatives, and specific sequences---in a
unified and efficient way.  The resulting automaton acts as the core
mechanism for scanning and detecting patterns within the data.

% ................................
\subsection{Execution Engine on Pandas}
\label{subsec:overview-execution}
% ................................

The execution phase uses the compiled DFA to perform pattern
matching directly on input data represented as Pandas DataFrames.
The engine begins by partitioning and ordering the data according to
the query's specifications.  It then applies the DFA to scan the
dataset, identifying all sequences that fit the defined pattern.
For every detected match the engine computes any requested measures
or calculations and compiles the results into a DataFrame formatted
to mirror standard SQL query outputs.  By utilising Pandas, the
system retains both flexibility and accessibility, smoothly
incorporating advanced pattern matching into conventional Python
data-analysis workflows.

% ----------------------------------------------------------------
\section{SQL Parsing}
\label{sec:meth-compiler}
% ----------------------------------------------------------------

Transforming a \texttt{MATCH\_RECOGNIZE} query from plain text into
an executable representation requires a parsing pipeline that
handles both standard SQL syntax and the extended row-pattern
sub-language.  This section explains parsing fundamentals, surveys
existing SQL parsers and parser generators, motivates the design
choices made in this implementation, and describes the resulting
parser in detail.

% ................................
\subsection{What is Parsing?}
\label{subsec:meth-what-parsing}
% ................................

\emph{Parsing} refers to the analysis of an input string according to the
rules of a formal grammar and its transformation into a
machine-processable representation~\cite{mogensen2024compiler}.  In essence,
the parser examines a human-written query, determines which grammar rules
govern its syntactic construction, and produces a tree data structure that
reflects that derivation.

Parsing in this engine proceeds in two sequential stages:
\emph{lexical analysis} and \emph{syntactic analysis}~\cite{DragonBook}.

\paragraph{Lexical analysis.}
The raw character stream is segmented into a sequence of predefined,
meaningful units called \emph{tokens}.  Keywords
(\texttt{SELECT}, \texttt{MATCH\_RECOGNIZE}, \texttt{PATTERN}, \ldots),
identifiers, string and numeric literals, and punctuation are all examples
of tokens.  Whitespace and comments are discarded at this stage; only the
ordered token sequence is passed to the next stage.

\paragraph{Syntactic analysis.}
The token stream is matched against the grammar rules of the language,
and a valid input is translated into a tree representation.  A valid
input produces a \emph{parse tree} (also called a \emph{concrete syntax
tree}) whose internal nodes correspond to grammar rules and whose leaves
correspond to individual tokens.  An invalid input causes the parser to
report an error, indicating precisely where the text violated the grammar.

The two-stage design separates concerns cleanly: the lexer handles
regular-language aspects (token shapes), while the parser handles
context-free-language aspects (phrase structure).  Together they convert
a character string into a structured tree that subsequent compiler phases
can traverse and transform.

% ................................
\subsection{Parsing SQL: Existing Solutions}
\label{subsec:meth-existing-parsers}
% ................................

Several Python libraries offer SQL parsing out of the box.
\texttt{sqlparse} performs tokenisation and basic structural
splitting (statements, clauses) but produces only a shallow token
stream, not a typed AST, and does not model \texttt{MATCH\_RECOGNIZE}
at all.  \texttt{sqlglot} is a more complete SQL transpiler that
supports multiple SQL dialects and exposes a typed AST; however,
its \texttt{MATCH\_RECOGNIZE} support is partial and does not cover
the full SQL:2016 pattern sub-language (e.g.\ exclusions, anchors,
\texttt{PERMUTE}, bounded quantifiers).  \texttt{mo\_sql\_parsing}
converts SQL to JSON but similarly provides no coverage for row
pattern constructs.

A widely used baseline for conventional SQL parsing is the
\textbf{PostgreSQL parser}, which has been extracted from the
PostgreSQL source code and released as a standalone C library via
the \texttt{libpg\_query} project.  Built on top of
\texttt{libpg\_query}, the Python library \texttt{pglast} wraps the
original C AST nodes as typed Python classes and provides
functionality to reconstruct SQL text from an AST.
Listing~\ref{lst:pglast-example} shows a minimal usage example:
the call to \texttt{parse\_sql} returns a tree of typed node
objects, and \texttt{RawStream} serialises it back to SQL.

\begin{lstlisting}[language=Python,
    caption={Parsing a SQL statement with \texttt{pglast} and
             reconstructing SQL text},
    label={lst:pglast-example}]
from pglast import parse_sql
from pglast.stream import RawStream

tree_root = parse_sql("SELECT 42;")
print(tree_root)          # typed AST node objects
print(RawStream()(tree_root))  # SELECT 42
\end{lstlisting}

Although \texttt{pglast} cannot parse \texttt{MATCH\_RECOGNIZE}, it
is useful for parsing the SQL sub-expressions that appear inside
\texttt{DEFINE} predicates and \texttt{MEASURES} clauses
(see Section~\ref{subsec:meth-safe-eval}).  The absence of any
library that handles the full SQL:2016 row-pattern sub-language
made a custom parser necessary for this project.

% ................................
\subsection{Grammars}
\label{subsec:meth-grammars}
% ................................

A grammar formally describes the syntax of a language~\cite{DragonBook}.
The elementary constituents of a grammar are:

\begin{itemize}
  \item \textbf{Terminal symbols} (tokens) --- the elementary,
    indivisible symbols of the language.  In this thesis the terms
    \emph{terminal} and \emph{token} are used interchangeably.
  \item \textbf{Non-terminals} --- symbols that represent sets of
    strings composed of terminals, used to name syntactic constructs
    (e.g.\ \emph{expression}, \emph{pattern}, \emph{clause}).
  \item \textbf{Production rules} --- rewriting rules of the form
    $N \to \alpha$, where $N$ is a non-terminal and $\alpha$ is a
    sequence of terminals and/or non-terminals.  Each rule describes
    one way in which a particular construct of the language may be
    written~\cite{wagenknecht2022formal}.
\end{itemize}

A grammar is specified as a list of production rules.  Because the
left-hand side of each rule contains exactly one non-terminal, the
grammar is \emph{context-free}: the applicable rules do not depend
on the surrounding context~\cite{DragonBook}.  For example, a grammar
for one-or-more occurrences of the letter $a$ can be written:

\[
  S \to a \quad\mid\quad aS
\]

This form is also known as \emph{Backus--Naur Form} (BNF).
The \emph{Extended} variant (EBNF) adds syntactic sugar for
optional elements (\texttt{[~]}), repetition (\texttt{\{~\}}), and
grouping (\texttt{(~)}), making large grammars more readable.  For
example, the \texttt{PARTITION~BY} clause can be expressed in EBNF
as:

\begin{center}
  \texttt{partitionBy~:~'PARTITION'~'BY'~expr~(','\ expr)*~;}
\end{center}

SQL is a large context-free language; the \texttt{MATCH\_RECOGNIZE}
sub-language adds several hundred additional productions for pattern
terms, quantifiers, alternation, grouping, anchors, exclusions,
\texttt{PERMUTE}, and the eight sub-clauses.

% ................................
\subsection{Our Parsing Needs}
\label{subsec:meth-parsing-needs}
% ................................

We now examine exactly what must be parsed and what the output
should look like.  Listing~\ref{lst:mr-full-syntax} shows the full syntax
of a query containing a \texttt{MATCH\_RECOGNIZE} clause as supported
by this implementation.

\begin{lstlisting}[
    language   = {},
    caption    = {Syntax of a \texttt{MATCH\_RECOGNIZE} query},
    label      = {lst:mr-full-syntax},
    basicstyle = \ttfamily\footnotesize
]
SELECT <expression> [, ...]
FROM   <table>
MATCH_RECOGNIZE (
  [ PARTITION BY <column> [, ...] ]
  [ ORDER BY <column> [ASC|DESC] [, ...] ]
  [ MEASURES <expression> [AS] <alias> [, ...] ]
  [ ONE ROW PER MATCH
  | ALL ROWS PER MATCH
    [ SHOW EMPTY MATCHES
    | OMIT EMPTY MATCHES
    | WITH UNMATCHED ROWS ] ]
  [ AFTER MATCH SKIP
    { PAST LAST ROW
    | TO NEXT ROW
    | TO FIRST <variable>
    | TO LAST  <variable> } ]
  PATTERN ( <row_pattern> )
  [ SUBSET <subset_name> = ( <variable> [, ...] ) [, ...] ]
  DEFINE <variable> AS <condition> [, ...]
) AS <alias>
<rest>;
\end{lstlisting}

\texttt{MATCH\_RECOGNIZE} is an optional sub-clause of the
\texttt{FROM} clause.  For simplicity, the implementation is
restricted to a single top-level \texttt{SELECT} statement:
no Common Table Expression (CTE), \texttt{VIEW} definition, or
inline sub-query nesting is supported.  Similarly, the trailing
\texttt{<rest>} part must not include set operations such as
\texttt{UNION}, \texttt{INTERSECT}, or \texttt{EXCEPT}.  These
restrictions are sufficient to cover the full SQL:2016 row-pattern
use-case within a Pandas-based execution engine.

The parser must extract all eight optional and mandatory
sub-clauses of \texttt{MATCH\_RECOGNIZE}---\texttt{PARTITION~BY},
\texttt{ORDER~BY}, \texttt{MEASURES}, rows-per-match,
after-match-skip, \texttt{PATTERN}, \texttt{SUBSET}, and
\texttt{DEFINE}---and deliver them to subsequent pipeline stages
through a well-defined AST interface, described next.

% ................................
\subsubsection{The \texttt{MATCH\_RECOGNIZE} AST Interface}
\label{subsubsec:meth-ast-interface}
% ................................

The AST exposed by the parser to downstream stages is defined by a
hierarchy of typed Python \texttt{@dataclass} objects.  The root
node, \texttt{FullQueryAST}, holds the outer \texttt{SelectClause},
\texttt{FromClause}, and a single \texttt{MatchRecognizeClause}
object that groups all eight sub-clause nodes.  The interface
contract is:

\begin{itemize}
  \item Every node is immutable after construction.
  \item Each \texttt{MATCH\_RECOGNIZE} sub-clause maps to a
    dedicated typed dataclass
    (\texttt{PartitionByClause}, \texttt{OrderByClause},
     \texttt{MeasuresClause}, \texttt{RowsPerMatchClause},
     \texttt{AfterMatchSkipClause}, \texttt{PatternClause},
     \texttt{SubsetClause}, \texttt{DefineClause}),
    enabling the NFA compiler to dispatch on type rather than
    parsing raw strings.
  \item \texttt{DEFINE} predicates and \texttt{MEASURES}
    expressions are captured as plain strings during the first
    pass and re-evaluated on demand through a lightweight
    expression evaluator (Section~\ref{subsec:meth-safe-eval}),
    keeping the main grammar concise.
  \item \texttt{Measure} nodes carry computed flags
    (\texttt{is\_classifier}, \texttt{is\_match\_number}) that
    allow the execution layer to route output columns without
    additional string inspection.
\end{itemize}

% ................................
\subsection{Parser Generators}
\label{subsec:meth-parser-generators}
% ................................

Since productivity is a goal, the question arises whether the
creation of a lexer and parser can be automated.  The kind of
software addressing this issue is called a
\emph{parser generator}~\cite{mogensen2024compiler}: it takes a
grammar as input and produces the source code of a parser that
understands the syntax defined by that grammar.  The target
language of the generated source code must be Python, since that
is the language the engine is written in.

Candidates that support Python as a target language include:

\begin{itemize}
  \item \textbf{PLY} (Python Lex-Yacc) --- a Python implementation
    of the classic 1970 Lex/Yacc tools~\cite{ply_tool}.
  \item \textbf{Lark} --- described as ``a modern approach, with a
    focus on ergonomics, performance and
    modularity''~\cite{lark_tool}.
  \item \textbf{ANTLR4} --- a parser generator ``widely used to
    build languages, tools, and frameworks; from a grammar, ANTLR
    generates a parser that can build and walk parse
    trees''~\cite{Parr2013}.
\end{itemize}

Table~\ref{tab:parser-gen-comparison} compares the three tools
across criteria extracted from their respective documentation.

\begin{table}[htbp]
  \centering
  \small
  \caption{Comparison of Python-compatible parser generators}
  \label{tab:parser-gen-comparison}
  \begin{tabular}{p{2.8cm} p{3.0cm} p{3.0cm} p{3.0cm}}
    \toprule
    \textbf{Feature}
      & \textbf{PLY (Lex/Yacc)}
      & \textbf{Lark}
      & \textbf{ANTLR4} \\
    \midrule
    Algorithm
      & LALR(1)
      & LALR / Earley
      & Adaptive LL(*) \cite{parr2014adaptive} \\[4pt]
    Grammar notation
      & BNF
      & EBNF
      & EBNF \\[4pt]
    Ambiguity handling
      & Operator precedence declarations
      & Rule priorities; Earley resolves via priority
      & Production order \& semantic predicates \\[4pt]
    Left recursion
      & Yes
      & Yes (LALR mode)
      & Yes \\[4pt]
    Target language
      & Python native
      & Python native
      & Java runtime required to generate Python parser \\[4pt]
    AST / parse tree
      & No
      & Yes (parse tree)
      & Yes (parse tree) \\[4pt]
    Tree visualiser
      & No
      & Dot output
      & Java GUI (\texttt{antlr4-tools}) \\[4pt]
    Traversal scaffolding
      & No
      & Transformer class; methods written manually
      & Listener \& visitor generated automatically \\
    \bottomrule
  \end{tabular}
\end{table}

Due to the manageable size and complexity of the
\texttt{MATCH\_RECOGNIZE} statement, any of the three tools would
likely produce a satisfactory result.  Nevertheless, two reasons
motivated the choice of ANTLR4.  First, its parse-tree visualiser
(via the \texttt{antlr4-tools} Java GUI) greatly aids debugging
grammar rules.  Second, and more crucially, ANTLR4 automatically
generates a traversal skeleton---both listener and visitor base
classes---for the parse tree.  Because the raw parse tree must be
transformed into the typed AST described in
Section~\ref{subsubsec:meth-ast-interface}, a generated visitor
base class eliminates manual tree-traversal boilerplate.  An
empirical study further demonstrates significantly higher
performance and maintainability when using ANTLR compared to
Lex/Yacc~\cite{antlr_performance}.

In this thesis the workflow combining ANTLR4 with the typed
\texttt{MATCH\_RECOGNIZE} AST interface is evaluated for its
suitability in parsing the full \texttt{MATCH\_RECOGNIZE} statement.

% ................................
\subsection{Alternative Methods}
\label{subsec:meth-alternatives}
% ................................

Two non-generator approaches were considered and rejected.

\paragraph{Hand-written recursive descent.}
A recursive descent parser is straightforward to implement for
small grammars: each non-terminal becomes a Python method that
calls other methods for sub-rules.  However, SQL's grammar contains
hundreds of productions; maintaining a hand-written parser as the
grammar evolves is error-prone.  The approach was ruled out on
maintainability grounds.

\paragraph{Regex-based string extraction.}
Regular expressions cannot model recursive structures such as nested
grouping \texttt{((A+B)*|C)} or balanced parentheses, so they
cannot serve as a complete parser.  They can pre-screen queries or
extract simple fragments, but any production use requires a proper
parser as the primary mechanism.

% ................................
\subsection{ANTLR4}
\label{subsec:meth-antlr4}
% ................................

\textbf{ANTLR4} (ANother Tool for Language Recognition, version~4)
is a parser generator that produces top-down, Adaptive LL(*)
recursive-descent parsers from context-free grammar specifications
written in the ANTLR meta-language~\cite{Parr2013,parr2014adaptive,antlr_website,parr_antlr_docs}.  Given a \texttt{.g4} grammar
file the tool generates a lexer, a parser, and listener/visitor
scaffolding in the target language (Python in this project).

The implementation extends Trino's production SQL grammar
(\texttt{TrinoLexer.g4} / \texttt{TrinoParser.g4}) with specialised
rules for every sub-clause of \texttt{MATCH\_RECOGNIZE}, as shown in
the excerpt in Listing~\ref{lst:grammar}.  A custom error listener
intercepts syntax errors and reports them with precise
line/column positions and contextual suggestions.

% ................................
\subsection{Implementation of the Parser}
\label{subsec:meth-parser-impl}
% ................................

% ................................
\subsubsection{Designing the Grammar}
\label{subsubsec:meth-grammar-design}
% ................................

The implementation uses the production Trino SQL grammar obtained
from the ANTLR community grammars project
(\texttt{TrinoLexer.g4} / \texttt{TrinoParser.g4},
approximately 1\,124 lines in total)~\cite{antlr_trino_sql}.
This grammar natively contains the \texttt{patternRecognition} rule
that covers all eight \texttt{MATCH\_RECOGNIZE} sub-clauses; no
modifications to the grammar file were required.
Listing~\ref{lst:grammar} shows the key productions.

\begin{lstlisting}[
    language = {},
    caption  = {Key productions of the \texttt{patternRecognition} rule in \texttt{TrinoParser.g4}},
    label    = {lst:grammar},
    basicstyle = \ttfamily\footnotesize
]
patternRecognition
    : aliasedRelation (
        MATCH_RECOGNIZE '('
          (PARTITION BY expression (',' expression)*)?
          (ORDER BY sortItem (',' sortItem)*)?
          (MEASURES measureDefinition (',' measureDefinition)*)?
          rowsPerMatch?
          (AFTER MATCH skipTo)?
          (INITIAL | SEEK)?
          PATTERN '(' rowPattern ')'
          (SUBSET subsetDefinition (',' subsetDefinition)*)?
          (DEFINE variableDefinition (',' variableDefinition)*)?
        ')' (AS? identifier columnAliases?)?
      )?
    ;

measureDefinition : expression AS identifier ;

rowsPerMatch
    : ONE ROW PER MATCH
    | ALL ROWS PER MATCH emptyMatchHandling?
    ;

skipTo
    : SKIP TO NEXT ROW
    | SKIP TO (FIRST | LAST)? identifier
    | SKIP PAST LAST ROW
    ;

rowPattern
    : rowPattern rowPattern             # patternConcatenation
    | rowPattern '|' rowPattern         # patternAlternation
    | patternPrimary patternQuantifier? # quantifiedPrimary
    ;

patternPrimary
    : identifier                                        # patternVariable
    | '(' rowPattern ')'                                # groupedPattern
    | '{-' rowPattern '-}'                              # excludedPattern
    | PERMUTE '(' rowPattern (',' rowPattern)* ')'      # patternPermutation
    ;
\end{lstlisting}

% ................................
\subsubsection{Handling Expressions as Strings}
\label{subsubsec:meth-expr-strings}
% ................................

\texttt{DEFINE} predicates and \texttt{MEASURES} expressions contain
arbitrary column references, navigation function calls, arithmetic,
and comparisons.  Parsing these into a fully typed expression tree
requires the grammar to enumerate every possible SQL operator and
function name.  An alternative---adopted here---is to capture such
expressions as structured token sequences during parsing, then
re-parse them through a lightweight expression parser on first use.
This two-pass strategy keeps the main grammar concise while still
producing a typed AST for safe evaluation (see
Section~\ref{subsec:meth-safe-eval}).

% ................................
\subsubsection{Parse Tree vs.\ Abstract Syntax Tree}
\label{subsubsec:meth-ast}
% ................................

Parse trees produced by ANTLR4 contain many structural nodes
(keywords, punctuation) that are unnecessary for execution.  A
visitor traverses the parse tree and produces a compact
\emph{Abstract Syntax Tree} (AST) whose node types correspond
directly to semantic constructs~\cite{DragonBook,jones2003ast}.  The key node
types are summarised in Table~\ref{tab:meth-ast-nodes}.

\begin{table}[htbp]
  \centering
  \caption{Principal AST node types}
  \label{tab:meth-ast-nodes}
  \begin{tabular}{ll}
    \toprule
    \textbf{Node Type}                  & \textbf{Represents} \\
    \midrule
    \texttt{FullQueryAST}               & Root of the full query (SELECT + FROM + MR) \\
    \texttt{MatchRecognizeClause}        & Root of the \texttt{MATCH\_RECOGNIZE} sub-clause \\
    \texttt{PartitionByClause}           & \texttt{PARTITION BY} column list \\
    \texttt{OrderByClause}              & \texttt{ORDER BY} sort items (\texttt{SortItem}) \\
    \texttt{MeasuresClause}             & \texttt{MEASURES} list (\texttt{Measure} nodes) \\
    \texttt{RowsPerMatchClause}          & \texttt{ONE / ALL ROWS PER MATCH} mode \\
    \texttt{AfterMatchSkipClause}        & \texttt{AFTER MATCH SKIP} policy \\
    \texttt{PatternClause}              & \texttt{PATTERN} string with tokenizer \\
    \texttt{SubsetClause}               & \texttt{SUBSET} variable definitions \\
    \texttt{DefineClause}               & \texttt{DEFINE} list (\texttt{Define} nodes) \\
    \bottomrule
  \end{tabular}
\end{table}

The AST is validated for semantic correctness before pattern
compilation: undefined variable references in \texttt{DEFINE}
predicates, circular definitions, and illegal quantifier ranges are
all detected and reported at this stage.

% ................................
\subsection{Visitor Implementation}
\label{subsec:meth-visitor}
% ................................

ANTLR4 generates an abstract visitor base class from the grammar~\cite{robson2022designpatterns}.
The parsing module is structured as two cooperating visitor classes
built on the generated \texttt{TrinoParserVisitor} base.

\paragraph{\texttt{MatchRecognizeExtractor}.}
The inner visitor overrides \texttt{visitPatternRecognition}
and extracts all eight sub-clauses in strict grammar order:
\texttt{PARTITION~BY} column expressions,
\texttt{ORDER~BY} sort items with direction and nulls handling,
\texttt{MEASURES} definitions (handling \texttt{RUNNING}/\texttt{FINAL}
prefix and alias via a dedicated \texttt{smart\_split()} helper that
correctly traverses nested parentheses and quoted identifiers),
rows-per-match mode, after-match-skip policy, the raw pattern string,
\texttt{SUBSET} definitions, and \texttt{DEFINE} variable--condition
pairs.  After extraction, a four-step validation chain is applied:
\texttt{validate\_clauses} (ensures \texttt{PATTERN}
and \texttt{DEFINE} are both present when required);
\texttt{validate\_pattern\_clause} (checks balanced parentheses and
flags conflicts such as pattern exclusion combined with
\texttt{WITH UNMATCHED ROWS});
\texttt{validate\_function\_usage} (permits only the whitelisted
navigation and aggregate functions---\texttt{FIRST}, \texttt{LAST},
\texttt{PREV}, \texttt{NEXT}, \texttt{COUNT}, \texttt{SUM\_IF},
\texttt{AVG\_IF}, \texttt{COUNT\_IF}---inside \texttt{MEASURES});
and \texttt{validate\_pattern\_variables\_defined} (verifies that
every variable referenced in the pattern appears in \texttt{DEFINE} or
\texttt{SUBSET}, and detects table-qualified references that would be
ambiguous at execution time).

\paragraph{\texttt{FullQueryExtractor}.}
The outer visitor handles the top-level \texttt{singleStatement}
rule.  It extracts the outer \texttt{SELECT} column list, \texttt{FROM}
table reference, and any trailing \texttt{ORDER~BY} clause via regex,
then locates the \texttt{PatternRecognitionContext} node through a
recursive tree walk and delegates to \texttt{MatchRecognizeExtractor}.

\paragraph{Public API.}
The entry point for all callers is
\texttt{parse\_full\_query(query:~str)} $\to$ \texttt{FullQueryAST}.
This function preprocesses the query string
(e.g.\ rewriting empty \texttt{IN~()} predicates to a sentinel value
to satisfy the grammar), constructs
\texttt{InputStream} $\to$ \texttt{TrinoLexer} $\to$ \texttt{CommonTokenStream} $\to$ \texttt{TrinoParser},
replaces the default ANTLR4 error listeners with a
\texttt{CustomErrorListener} that raises a structured
\texttt{ParserError} carrying exact line, column, and source snippet
on any syntax violation, and visits the parse tree with
\texttt{FullQueryExtractor} to produce the typed \texttt{FullQueryAST}.

% ................................
\subsection{Testing and Evaluation of the Parser}
\label{subsec:meth-parser-testing}
% ................................

The parser implementation is validated through a structured test
suite located in the \texttt{tests/} directory.  Tests are organised
into three layers: unit tests for individual extractor methods,
integration tests that exercise the full
\texttt{parse\_full\_query} pipeline, and
SQL:2016 compliance tests that compare output against Trino's
reference behaviour.

\paragraph{Unit tests.}
Each of the eight \texttt{extract\_*} methods is tested in isolation
against known SQL fragments.  For example, \texttt{extract\_measures}
is fed a \texttt{MEASURES} token stream containing \texttt{RUNNING},
\texttt{FINAL}, and undecorated expressions, and the resulting
\texttt{Measure} objects are asserted for correct alias, expression
text, and semantic flags (\texttt{explicit\_semantics},
\texttt{is\_classifier}, \texttt{is\_match\_number}).
The validation chain is also tested adversarially: supplying an
undefined pattern variable, a prohibited function in \texttt{MEASURES},
or a \texttt{PATTERN} clause without a corresponding \texttt{DEFINE}
all trigger a \texttt{ParserError} carrying the expected
line and column information.

\paragraph{Integration tests.}
The file \texttt{test\_match\_recognize.py} mirrors Trino's
\texttt{TestRowPatternMatching.java} and exercises twelve scenarios
end-to-end: simple queries, pattern quantifiers (\texttt{*}, \texttt{+},
\texttt{?}, bounded \texttt{\{n,m\}}), pattern exclusion syntax
(\texttt{\{-\ldots-\}}), all four after-match-skip modes,
one-row and all-rows output modes, \texttt{CLASSIFIER()} and
\texttt{MATCH\_NUMBER()} special columns, all four navigation functions,
partitioned and ordered input, \texttt{SUBSET} aliasing, and
\texttt{RUNNING}/\texttt{FINAL} semantics.

\paragraph{SQL:2016 compliance tests.}
\texttt{test\_sql2016\_compliance.py} contains seven test classes
targeted specifically at standard conformance: pattern syntax,
navigation functions, special aggregate functions, partition
boundary handling, subset variable resolution, output-mode
semantics, and all six skip policies.  Empty-match handling is
covered separately in \texttt{test\_empty\_matches.py} with nine
test cases including \texttt{SHOW EMPTY MATCHES},
\texttt{OMIT EMPTY MATCHES}, \texttt{WITH UNMATCHED ROWS}, and
combinations with alternation and partitioning.

\paragraph{Edge-case tests.}
Additional test files cover \texttt{IN~()} predicate preprocessing
(\texttt{test\_in\_predicate.py}),
\texttt{IS~NULL} handling (\texttt{test\_is\_null\_handling.py}),
anchored patterns (\texttt{test\_anchor\_patterns.py}),
\texttt{PERMUTE} patterns (\texttt{test\_permute\_patterns.py}),
back-references (\texttt{test\_back\_reference.py}),
pattern cache correctness and thread safety
(\texttt{test\_pattern\_cache.py}), and
production-scale aggregation accuracy
(\texttt{test\_production\_aggregates.py}).

% ................................
\subsection{Discussion of the Parser}
\label{subsec:meth-parser-discussion}
% ................................

The grammar-driven approach provides several practical advantages:
the language specification remains in a single readable
\texttt{.g4} file; ANTLR4's error recovery inserts synthetic tokens
to continue parsing past a local error, collecting multiple
diagnostics in one pass; and the generated visitor scaffolding
eliminates manual tree traversal boilerplate.

% ................................
\subsubsection{Limits}
\label{subsubsec:meth-parser-limits}
% ................................

Two limitations are known.  First, some advanced SQL:2016 expression
constructs (e.g.\ recursive \texttt{WITH RECURSIVE} inside a
\texttt{DEFINE} predicate) are not supported; the expression
sub-parser recognises only the constructs needed for
\texttt{MATCH\_RECOGNIZE} evaluation.  Second, because the grammar
extends Trino's dialect, vendor-specific extensions of other
systems (Oracle \texttt{RUNNING} keyword placement differences,
Snowflake aggregate variants) require minor grammar adjustments if
cross-dialect compatibility is needed.

% ................................
\subsection{Safe SQL Expression Evaluation}
\label{subsec:meth-safe-eval}
% ................................

\texttt{DEFINE} predicates and \texttt{MEASURES} expressions are
user-supplied strings that must be evaluated as Python code at
runtime.  Directly calling \texttt{eval()} on arbitrary user input
creates a code-injection vulnerability.  The implementation instead
uses Python's \texttt{ast} module to parse and validate expressions
before execution, as shown in Listing~\ref{lst:safe-eval}.

\begin{lstlisting}[
    language = Python,
    caption  = {Safe SQL expression evaluation using AST whitelisting},
    label    = {lst:safe-eval}
]
import ast

def safe_evaluate(sql_expr, row_data, context):
    # 1. Translate SQL syntax to Python equivalents
    py_expr = translate_sql_to_python(sql_expr)

    # 2. Parse into a Python AST (expression mode only)
    tree = ast.parse(py_expr, mode='eval')

    # 3. Whitelist-validate: allow comparisons, arithmetic,
    #    boolean ops, literals, subscripts -- forbid imports,
    #    function calls (except navigation functions), assignments
    ASTSecurityValidator().visit(tree)

    # 4. Compile and execute in a controlled namespace
    code = compile(tree, '<expr>', 'eval')
    return eval(code, {'__builtins__': {}},
                       {'row': row_data, **context})
\end{lstlisting}

The controlled namespace exposes only the current row dictionary and
the whitelisted navigation functions (\texttt{PREV}, \texttt{NEXT},
\texttt{FIRST}, \texttt{LAST}), providing a safe sandbox with no
access to the file system, network, or interpreter internals.


% ----------------------------------------------------------------
\section{Semantic Validation}
\label{sec:meth-semantic}
% ----------------------------------------------------------------

After the AST is constructed, the engine performs a systematic
semantic validation pass before any automaton is built.  This pass
catches errors that are syntactically valid but logically
inconsistent or non-conformant with the SQL:2016 standard, and
applies a set of normalisation rules that simplify downstream
processing.

% ................................
\subsection{Schema and Variable Consistency Checks}
\label{subsec:sem-schema}
% ................................

The validator first verifies that:

\begin{itemize}
  \item Every column name referenced in \texttt{PARTITION BY},
    \texttt{ORDER BY}, \texttt{DEFINE}, and \texttt{MEASURES}
    expressions exists in the input DataFrame schema.
    Column names are resolved case-insensitively to match
    common SQL conventions.

  \item Every pattern variable referenced in \texttt{MEASURES} or
    \texttt{DEFINE} expressions has been declared in the
    \texttt{PATTERN} clause.  Variables appearing in
    \texttt{DEFINE} but absent from \texttt{PATTERN} are treated
    as an error rather than silently ignored.

  \item If a \texttt{SUBSET} clause is present, each component
    variable listed in the subset definition must be a declared
    pattern variable.  Circular subset definitions are rejected.

  \item The \texttt{AFTER MATCH SKIP TO FIRST/LAST $V$} options
    reference a pattern variable $V$ that actually exists in the
    pattern, catching typos at compile time rather than at runtime.
\end{itemize}

% ................................
\subsection{Navigation-Function Validation}
\label{subsec:sem-navfunc}
% ................................

Navigation functions impose inter-variable ordering constraints that
must be verified before the automaton is constructed.  The
\texttt{validate\_navigation\_conditions()} function in
\texttt{src/executor/match\_recognize.py} performs this check by
building a \emph{variable timeline}: the order in which distinct
pattern variables first appear in the pattern string.  For each
\texttt{DEFINE} predicate the validator then checks:

\begin{itemize}
  \item \texttt{NEXT($v$.$c$)} is only permitted inside the
    predicate of a non-final pattern variable, i.e.\ a variable
    that is not last in the timeline.  Using \texttt{NEXT} in the
    predicate of the final variable would always reference a row
    beyond the match, violating SQL:2016.  If this rule is
    violated, the engine returns an empty result set rather than
    raising an exception---mirroring the behaviour documented for
    Trino.

  \item \texttt{PREV($v$.$c$, $k$)} references are accepted in any
    variable's predicate; the validator checks that the offset $k$
    is a non-negative integer literal.

  \item Nested navigation expressions (e.g.\
    \texttt{NEXT(PREV(A.price, 1), 1)}) are validated recursively
    by \texttt{validate\_complex\_nested\_navigation()}.
\end{itemize}

% ................................
\subsection{Output-Mode and Skip-Mode Consistency}
\label{subsec:sem-modes}
% ................................

Several combinations of output mode, skip mode, and pattern
features are prohibited by the SQL:2016 standard.  The validator
enforces the following rules:

\begin{itemize}
  \item Pattern exclusion syntax (\texttt{\{- \ldots -\}}) is
    incompatible with \texttt{ALL ROWS PER MATCH WITH UNMATCHED
    ROWS}; this combination is rejected with a descriptive error
    message.

  \item \texttt{AFTER MATCH SKIP TO FIRST/LAST $V$} requires that
    $V$ be a pattern variable, not a subset alias.  Subset aliases
    are resolved at automaton-construction time, not at skip-mode
    computation time, so this restriction is necessary to guarantee
    a deterministic next-start position.
\end{itemize}

\noindent
After all checks pass, the validator records the resolved
configuration—normalised output mode, skip mode, implicit-variable
defaults, and measure semantics—ready for the logical planning step.

% ................................................................
\section{Logical Plan Construction}
\label{sec:meth-logical}
% ................................................................

The logical plan translates the validated AST into a flat,
canonical data structure that drives all subsequent phases without
further reference to the raw SQL text.  In the implementation this
materialises as a set of typed Python objects extracted from the
\texttt{MatchRecognizeClause} AST node and passed as explicit
arguments to the execution functions.
Table~\ref{tab:logical-plan} lists the principal logical-plan
attributes and their types.

\begin{table}[htbp]
  \centering
  \caption{Principal logical-plan attributes extracted from the
    validated AST}
  \label{tab:logical-plan}
  \begin{tabular}{llp{6.5cm}}
    \toprule
    \textbf{Attribute} & \textbf{Type} & \textbf{Semantics} \\
    \midrule
    \texttt{partition\_by}   & \texttt{List[str]}
      & Column names for DataFrame partitioning; empty list means
        single global partition \\[4pt]
    \texttt{order\_by}       & \texttt{List[str]}
      & Column names for intra-partition sort order \\[4pt]
    \texttt{pattern\_text}   & \texttt{str}
      & Normalised pattern string forwarded to the tokeniser \\[4pt]
    \texttt{define}          & \texttt{Dict[str, str]}
      & Variable $\to$ predicate expression mapping; implicit
        variables (absent from \texttt{DEFINE}) are assigned
        the always-true predicate \texttt{"True"} \\[4pt]
    \texttt{subset\_dict}    & \texttt{Dict[str, List[str]]}
      & Subset alias $\to$ constituent variable list, extracted
        by \texttt{extract\_subset\_dict()} \\[4pt]
    \texttt{measures}        & \texttt{Dict[str, str]}
      & Measure alias $\to$ expression string \\[4pt]
    \texttt{measure\_semantics} & \texttt{Dict[str, str]}
      & Measure alias $\to$ \texttt{"RUNNING"} or
        \texttt{"FINAL"}, determined per the rules summarised in
        Table~\ref{tab:rows-per-match} \\[4pt]
    \texttt{rows\_per\_match} & \texttt{RowsPerMatch}
      & Output-row multiplicity enum value \\[4pt]
    \texttt{skip\_mode}      & \texttt{SkipMode}
      & After-match skip strategy enum value \\[4pt]
    \texttt{skip\_var}       & \texttt{Optional[str]}
      & Target variable for \textsc{To First}/\textsc{To Last};
        \texttt{None} for the other two modes \\
    \bottomrule
  \end{tabular}
\end{table}

% ................................
\subsection{Implicit Variable Defaults}
\label{subsec:logical-implicit}
% ................................

SQL:2016 permits pattern variables to be listed in the
\texttt{PATTERN} clause without a corresponding entry in the
\texttt{DEFINE} clause.  Such variables match any row
unconditionally.  During logical plan construction the engine
identifies these implicit variables by comparing the set of names
extracted from the pattern tokens against the keys of the
\texttt{define} dictionary.  Any variable present in the former
but absent from the latter is assigned the always-true predicate
(represented as the Python string \texttt{"True"}), so that later
stages can treat all variables uniformly without special-casing
implicit ones.

% ................................
\subsection{Measure Semantics Resolution}
\label{subsec:logical-measure-sem}
% ................................

Measure semantics (\textsc{Running} vs.\ \textsc{Final}) can be
specified explicitly using the keywords \texttt{RUNNING} and
\texttt{FINAL} in the \texttt{MEASURES} clause, or left implicit
for the engine to determine.  The resolution algorithm
(Figure~\ref{fig:measure-sem-resolution}) scans all measure
expressions in order:

\begin{enumerate}
  \item If any measure carries an explicit \texttt{RUNNING} or
    \texttt{FINAL} keyword, the query is in \emph{mixed-semantics
    mode} and all other measures without an explicit keyword
    default to \textsc{Final}.

  \item In single-semantics mode (no explicit keywords):
    \begin{itemize}
      \item All measures default to \textsc{Final} under
        \texttt{ONE ROW PER MATCH}.
      \item Under \texttt{ALL ROWS PER MATCH}: \texttt{LAST()}
        defaults to \textsc{Running}; \texttt{FIRST()},
        \texttt{PREV()}, and \texttt{NEXT()} default to
        \textsc{Final}; aggregate functions \texttt{COUNT()},
        \texttt{SUM()}, and \texttt{ARRAY\_AGG()} default to
        \textsc{Running}; all others default to \textsc{Final}.
    \end{itemize}
\end{enumerate}

\begin{figure}[htbp]
  \centering
  \begin{tikzpicture}[
    node distance=0.5cm and 1.8cm,
    every node/.style={font=\small},
    decision/.style={diamond, draw, aspect=2.2,
                     inner sep=1pt, align=center},
    result/.style={rectangle, draw, rounded corners=3pt,
                   minimum width=2.8cm, minimum height=0.6cm,
                   align=center},
    arrow/.style={->, thick}
  ]
    \node[decision] (explicit) {Explicit\\keyword?};
    \node[result, right=2.2cm of explicit] (useexp) {Use stated\\semantics};
    \node[decision, below=1.2cm of explicit] (mixed) {Any measure\\explicit?};
    \node[result, right=2.2cm of mixed] (final0) {\textsc{Final}};
    \node[decision, below=1.2cm of mixed] (onerow) {\texttt{ONE ROW}\\mode?};
    \node[result, right=2.2cm of onerow] (final1) {\textsc{Final}};
    \node[decision, below=1.2cm of onerow] (functype) {\texttt{LAST}/\texttt{COUNT}/\\
      \texttt{SUM}/\texttt{ARRAY\_AGG}?};
    \node[result, right=2.2cm of functype] (running) {\textsc{Running}};
    \node[result, below=0.9cm of functype] (final2) {\textsc{Final}};

    \draw[arrow] (explicit) -- node[above,font=\tiny]{yes} (useexp);
    \draw[arrow] (explicit) -- node[left,font=\tiny]{no} (mixed);
    \draw[arrow] (mixed) -- node[above,font=\tiny]{yes} (final0);
    \draw[arrow] (mixed) -- node[left,font=\tiny]{no} (onerow);
    \draw[arrow] (onerow) -- node[above,font=\tiny]{yes} (final1);
    \draw[arrow] (onerow) -- node[left,font=\tiny]{no} (functype);
    \draw[arrow] (functype) -- node[above,font=\tiny]{yes} (running);
    \draw[arrow] (functype) -- node[left,font=\tiny]{no} (final2);
  \end{tikzpicture}
  \caption{Measure semantics resolution algorithm applied to each
    measure expression individually.  The decision tree follows
    SQL:2016 default-semantics rules.}
  \label{fig:measure-sem-resolution}
\end{figure}

\noindent
Once all attributes are resolved, the logical plan is complete and
is passed to the pattern-compilation phase described in the next
section.

% ----------------------------------------------------------------
\section{Theoretical Foundations: Finite Automata}
\label{sec:meth-automata}
% ----------------------------------------------------------------

Pattern matching in \texttt{MATCH\_RECOGNIZE} is fundamentally a
problem of sequence recognition.  A sequence of data rows, each
labelled with the pattern variable(s) it satisfies, constitutes a
string over the alphabet of pattern variables.  The problem of
determining whether this string is accepted by the pattern is
exactly the problem solved by finite automata~\cite{HopcroftUllman}.

% ................................
\subsection{Non-deterministic Finite Automata (NFA)}
\label{subsec:meth-nfa}
% ................................

A \emph{Non-deterministic Finite Automaton} (NFA) is formally
defined as a 5-tuple $M = (Q,\,\Sigma,\,\delta,\,q_0,\,F)$, where:

\begin{itemize}
  \item $Q$ is a finite set of \emph{states};
  \item $\Sigma$ is the \emph{input alphabet} (pattern variables in
    our setting);
  \item $\delta : Q \times (\Sigma \cup \{\varepsilon\}) \to 2^{Q}$
    is the \emph{transition function}, mapping each state and symbol
    (or the empty string $\varepsilon$) to a \emph{set} of successor
    states;
  \item $q_0 \in Q$ is the \emph{initial state}; and
  \item $F \subseteq Q$ is the set of \emph{accepting states}.
\end{itemize}

NFAs naturally model alternation and optional constructs owing to
their ability to be in multiple states simultaneously and to use
$\varepsilon$-transitions (transitions that consume no input).

Figure~\ref{fig:meth-nfa} shows the NFA for pattern $A^{+}\,B?$.

\begin{figure}[htbp]
  \centering
  \begin{tikzpicture}[
      >=stealth, shorten >=1pt,
      node distance=2.8cm, auto, thick,
      every state/.style={circle, draw, minimum size=1cm}
    ]
    \node[state, initial, initial text={}]  (q0) {$q_0$};
    \node[state]                            (q1) [right of=q0] {$q_1$};
    \node[state]                            (q2) [right of=q1] {$q_2$};
    \node[state, accepting]                 (q3) [right of=q2] {$q_3$};

    \path[->]
      (q0) edge              node {$A$}           (q1)
      (q1) edge [loop above] node {$A$}           (q1)
      (q1) edge              node {$\varepsilon$} (q2)
      (q2) edge              node {$B$}           (q3)
      (q2) edge [bend right=35] node [below] {$\varepsilon$} (q3);
  \end{tikzpicture}
  \caption{NFA for the pattern $A^{+}\,B?$ illustrating
    $\varepsilon$-transitions and non-determinism}
  \label{fig:meth-nfa}
\end{figure}

% ................................
\subsection{Thompson's Construction}
\label{subsec:meth-thompson}
% ................................

Thompson's construction algorithm~\cite{Thompson1968} converts a
regular expression into an NFA in linear time $O(n)$, where $n$ is
the length of the expression.  The algorithm proceeds recursively
over the expression structure according to the rules in
Algorithm~\ref{alg:thompson}.

\begin{algorithm}[htbp]
\caption{Thompson's NFA Construction}
\label{alg:thompson}
\DontPrintSemicolon
\SetAlgoLined
\SetKwProg{Proc}{Procedure}{:}{}
\SetKw{KwRet}{return}
\Proc{\upshape\textsc{Build}$(e)$}{
  \uIf{$e = \varepsilon$}{
    \KwRet NFA with one $\varepsilon$-transition from start to accept\;
  }
  \uElseIf{$e = a \in \Sigma$}{
    \KwRet NFA with one $a$-labelled transition from start to accept\;
  }
  \uElseIf{$e = r \cdot s$ \textnormal{(concatenation)}}{
    $N_r \leftarrow \textsc{Build}(r)$\;
    $N_s \leftarrow \textsc{Build}(s)$\;
    Merge accept state of $N_r$ with start state of $N_s$\;
    \KwRet combined NFA\;
  }
  \uElseIf{$e = r \mid s$ \textnormal{(alternation)}}{
    $N_r \leftarrow \textsc{Build}(r)$\;
    $N_s \leftarrow \textsc{Build}(s)$\;
    Add new start $q_\text{new}$ with $\varepsilon$-edges to $q_0^{N_r}$ and $q_0^{N_s}$\;
    Add new accept $f_\text{new}$ with $\varepsilon$-edges from $F^{N_r}$ and $F^{N_s}$\;
    \KwRet combined NFA\;
  }
  \uElseIf{$e = r^{*}$ \textnormal{(Kleene star)}}{
    $N_r \leftarrow \textsc{Build}(r)$\;
    Add $\varepsilon$-edge: start $\to$ accept (zero repetitions)\;
    Add $\varepsilon$-edge: $F^{N_r} \to q_0^{N_r}$ (loop back)\;
    \KwRet modified NFA\;
  }
  \ElseIf{$e = r^{\{n,m\}}$ \textnormal{(bounded quantifier)}}{
    Construct $n$ mandatory copies followed by $(m{-}n)$ optional copies\;
    \KwRet combined NFA\;
  }
}
\end{algorithm}

The resulting NFA has at most $2n$ states and is constructed in
$O(n)$ time and space, making it suitable for patterns of arbitrary
complexity~\cite{DragonBook}.

% ................................
\subsection{Pattern Tokenisation}
\label{subsec:meth-pattern-tokeniser}
% ................................

Before an NFA can be constructed the raw pattern string extracted
from the SQL query must be converted into a structured, traversable
representation.  The \texttt{tokenize\_pattern()} function in
\texttt{src/matcher/pattern\_tokenizer.py} performs this step in a
single left-to-right scan, producing a flat list of
\texttt{PatternToken} objects—the \emph{linearised} form of the
pattern's syntax tree.

\paragraph{PatternToken structure.}
Each \texttt{PatternToken} is a Python \texttt{dataclass} with the
following key fields:

\begin{itemize}
  \item \textbf{\texttt{type}} — a \texttt{PatternTokenType} enum
    value determining which NFA fragment builder is invoked
    (see Table~\ref{tab:token-types});
  \item \textbf{\texttt{value}} — the raw textual value extracted
    from the pattern string (variable name, operator character, etc.);
  \item \textbf{\texttt{quantifier}} — optional quantifier string
    (\texttt{'*'}, \texttt{'+'}, \texttt{'?'}, or
    \texttt{'\{$n$,$m$\}'}) copied from the immediately following
    quantifier when present; defaults to \texttt{None};
  \item \textbf{\texttt{greedy}} — boolean flag indicating greedy
    (\texttt{True}, the default) or reluctant (\texttt{False})
    quantification, set when a trailing \texttt{?} follows the
    quantifier;
  \item \textbf{\texttt{metadata}} — a dictionary carrying
    type-specific structured data (e.g.\ a \texttt{variables} list
    for \texttt{PERMUTE} tokens, an \texttt{excluded\_variables}
    list for \texttt{EXCLUSION} tokens); and
  \item \textbf{\texttt{priority}} — an integer priority inherited
    from the token's position in the stream, propagated to NFA
    transitions to enforce greedy/leftmost semantics.
\end{itemize}

\begin{table}[htbp]
  \centering
  \caption{The thirteen \texttt{PatternTokenType} values and their
    roles in the pattern string}
  \label{tab:token-types}
  \begin{tabular}{lp{8cm}}
    \toprule
    \textbf{Token type}         & \textbf{Role in pattern} \\
    \midrule
    \texttt{LITERAL}            & A pattern variable name (\texttt{A}, \texttt{UP},
                                  \texttt{STRT}, etc.) \\
    \texttt{ALTERNATION}        & The \texttt{|} operator \\
    \texttt{PERMUTE}            & A complete \texttt{PERMUTE(\ldots)} construct
                                  carrying its variable list in \texttt{metadata} \\
    \texttt{GROUP\_START}       & An opening parenthesis \texttt{(} \\
    \texttt{GROUP\_END}         & A closing parenthesis \texttt{)} \\
    \texttt{ANCHOR\_START}      & The start-of-partition anchor \texttt{\^{}} \\
    \texttt{ANCHOR\_END}        & The end-of-partition anchor \texttt{\$} \\
    \texttt{EXCLUSION}          & A complete \texttt{\{- \ldots -\}} exclusion block \\
    \texttt{EXCLUSION\_START}   & Opening delimiter \texttt{\{-} \\
    \texttt{EXCLUSION\_END}     & Closing delimiter \texttt{-\}} \\
    \texttt{QUANTIFIER}         & A standalone quantifier node \\
    \texttt{SUBSET}             & A \texttt{SUBSET} variable reference \\
    \texttt{SEQUENCE}           & A composite sequence node \\
    \bottomrule
  \end{tabular}
\end{table}

\paragraph{Tokenisation algorithm.}
The internal function \texttt{\_tokenize\_pattern\_internal()} scans
the pattern string character by character, dispatching on the current
position:

\begin{enumerate}
  \item Whitespace is skipped unconditionally.
  \item The keyword \texttt{PERMUTE} (case-insensitive) triggers
    \texttt{\_parse\_permute\_pattern()}, which collects the
    comma-separated variable list inside the following parentheses
    and returns a single \texttt{PERMUTE} token whose
    \texttt{metadata[`variables']} holds the parsed variable tokens.
  \item The two-character sequence \texttt{\{-} triggers
    \texttt{\_parse\_exclusion\_pattern()}, which reads to the
    matching \texttt{-\}} (tracking nesting depth) and returns an
    \texttt{EXCLUSION} token.
  \item The structural characters \texttt{(}, \texttt{)},
    \texttt{|}, \texttt{\^{}}, and \texttt{\$} each produce a
    single-character token of the corresponding type.
  \item Quantifier characters \texttt{*}, \texttt{+}, \texttt{?},
    and \texttt{\{} are parsed by \texttt{parse\_quantifier\_at()}
    and \emph{attached} to the preceding token as its
    \texttt{quantifier} and \texttt{greedy} fields—no separate
    token is inserted.
  \item Identifier characters (\texttt{[A-Za-z0-9\_]}) are
    accumulated into a \texttt{LITERAL} token.
\end{enumerate}

After the scan, \texttt{\_optimize\_token\_sequence()} makes a
conservative pass that removes structural redundancies while
preserving the one-variable-per-token invariant that the NFA builder
requires: adjacent \texttt{LITERAL} tokens that look like pattern
variable names are deliberately \emph{not} merged.

\paragraph{Quantifier normalisation.}
Quantifiers are stored as attributes of the token they modify rather
than as separate tokens.  The helper \texttt{parse\_quantifier()}
normalises all surface forms to the triple
$(\textit{min},\,\textit{max},\,\textit{greedy})$:
\begin{center}
\begin{tabular}{lll}
  \texttt{?} & $\to$ & $(0,\;1,\;\mathit{True})$ \\
  \texttt{*} & $\to$ & $(0,\;\infty,\;\mathit{True})$ \\
  \texttt{+} & $\to$ & $(1,\;\infty,\;\mathit{True})$ \\
  \texttt{\{$n$\}} & $\to$ & $(n,\;n,\;\mathit{True})$ \\
  \texttt{\{$n$,$m$\}} & $\to$ & $(n,\;m,\;\mathit{True})$ \\
  append \texttt{?} to any & $\to$ & $\mathit{greedy} = \mathit{False}$
\end{tabular}
\end{center}

\noindent This triple is consumed directly by
\texttt{NFABuilder.\_apply\_quantifier()} during NFA construction.

% ................................
\subsection{NFA Construction from Tokens}
\label{subsec:meth-nfa-from-tokens}
% ................................

The \texttt{NFABuilder} class in \texttt{src/matcher/automata.py}
implements the \emph{visitor} pattern over the token list.  Its
central method \texttt{\_process\_sequence()} iterates the token
stream and dispatches on each token's \texttt{type} field to invoke
the appropriate NFA-fragment builder—mirroring the recursive
tree-traversal structure of classical Thompson's construction but
applied to the flat linearised token representation rather than an
explicit parse tree.

\paragraph{Overall build sequence.}
The public entry point \texttt{NFABuilder.build()} executes the
following steps:

\begin{enumerate}
  \item Check the LRU pattern cache (keyed on an MD5 hash of the
    token list and DEFINE context) for a previously compiled NFA.
    A cache hit returns a deep clone of the stored NFA to prevent
    state sharing.
  \item On a cache miss call \texttt{\_build\_nfa\_internal()}, which:
    \begin{enumerate}[label=(\alph*)]
      \item validates anchor semantics (e.g.\ rejecting patterns such
        as \texttt{A\^{}} where a start anchor follows a literal);
      \item creates a global start state $q_\text{start}$ and accept
        state $q_\text{accept}$;
      \item delegates to \texttt{\_process\_sequence()} for the
        pattern body; and
      \item connects the pattern fragment to the global states via
        $\varepsilon$-transitions.
    \end{enumerate}
  \item Apply \texttt{NFA.optimize()} to remove unreachable states,
    deduplicate transitions, and pre-compute DFA hot-paths.
  \item Cache the compiled NFA and return it.
\end{enumerate}

\subsubsection{Token-to-NFA Fragment Mapping}
\label{subsubsec:meth-token-to-nfa}

Table~\ref{tab:token-nfa-map} summarises how the visitor translates
each token type into NFA structure.

\begin{table}[htbp]
  \centering
  \caption{Token types and their NFA fragment builders in
    \texttt{NFABuilder.\_process\_sequence()}}
  \label{tab:token-nfa-map}
  \begin{tabular}{lp{8.2cm}}
    \toprule
    \textbf{Token type}            & \textbf{NFA construction rule} \\
    \midrule
    \texttt{LITERAL}               & \texttt{create\_var\_states()} produces a
                                     two-state fragment
                                     $q_s \xrightarrow{C(v)} q_e$, where
                                     $C(v)$ is the callable compiled from the
                                     \texttt{DEFINE} predicate for variable $v$;
                                     the token's \texttt{quantifier} (if any) is
                                     applied by \texttt{\_apply\_quantifier()} \\[4pt]
    \texttt{ALTERNATION}           & The left branch is finalised; the right branch
                                     is processed by a recursive call to
                                     \texttt{\_process\_sequence()}; a fresh start
                                     state fans out with $\varepsilon$-transitions
                                     to both branch heads, and a fresh end state
                                     collects both branch tails \\[4pt]
    \texttt{GROUP\_START}          & Triggers a recursive call to
                                     \texttt{\_process\_sequence()} that reads
                                     until the matching \texttt{GROUP\_END}; the
                                     group quantifier (if any) is applied to the
                                     resulting fragment \\[4pt]
    \texttt{PERMUTE}               & \texttt{\_process\_permute()} expands all
                                     $n!$ orderings of the $n$ variables and
                                     builds a union of $n!$ concatenation
                                     fragments under a shared start/end pair \\[4pt]
    \texttt{EXCLUSION}             & \texttt{\_process\_exclusion\_token()} builds
                                     the contained sub-pattern normally but marks
                                     every state in the range with
                                     \texttt{is\_excluded = True}; the matcher
                                     filters excluded variables from output at
                                     match time \\[4pt]
    \texttt{ANCHOR\_START}/\texttt{END}
                                   & Creates a single anchor state with
                                     \texttt{is\_anchor = True} and the
                                     corresponding \texttt{anchor\_type};
                                     partition-boundary semantics are enforced
                                     by the matcher \\
    \bottomrule
  \end{tabular}
\end{table}

\subsubsection{Quantifier Expansion}
\label{subsubsec:meth-apply-quantifier}

\texttt{\_apply\_quantifier(start, end, }$\mathit{min}$\texttt{, }$\mathit{max}$\texttt{, }$\mathit{greedy}$\texttt{)}
wraps a fragment $(s,\,e)$ in outer envelope states $(q_s,\,q_e)$
and adds $\varepsilon$-transitions accordingly:

\begin{itemize}
  \item \textbf{\texttt{?}} $(\min=0,\,\max=1)$: adds
    $\varepsilon(q_s \to s)$ for the match case and
    $\varepsilon(q_s \to q_e)$ for the skip case.
  \item \textbf{\texttt{*}} $(\min=0,\,\max=\infty)$: adds
    $\varepsilon(q_s \to s)$, $\varepsilon(e \to q_e)$ (exit),
    $\varepsilon(q_s \to q_e)$ (skip), and
    $\varepsilon(e \to s)$ (loop for repetition).
  \item \textbf{\texttt{+}} $(\min=1,\,\max=\infty)$: as \texttt{*}
    but without the skip transition, enforcing at least one match.
  \item \textbf{\texttt{\{$n$,$m$\}}}: constructs $n$ mandatory
    copies of the fragment in series (connected by
    $\varepsilon$-transitions), followed by $(m-n)$ optional copies
    each of which carries an $\varepsilon$-bypass to $q_e$—the
    standard elimination-of-quantifier reduction.
  \item \textbf{Reluctant} (\texttt{greedy = False}): the same
    $\varepsilon$ structure is used, but the loop-back and
    extension transitions are annotated with a \emph{higher}
    numeric priority value so the epsilon-closure traversal
    prefers the exit path, yielding shorter matches.
\end{itemize}

\subsubsection{PERMUTE Expansion}
\label{subsubsec:meth-permute-expansion}

The \texttt{PERMUTE}$(v_1, v_2, \ldots, v_n)$ construct is
semantically equivalent to the alternation of all $n!$ orderings of
its arguments.  \texttt{\_process\_permute()} proceeds in three
phases:

\begin{enumerate}
  \item \textbf{Complexity guard.}  Patterns with more than eight
    variables are rejected immediately (raising
    \texttt{RuntimeError}) because $8! = 40\,320$ is the largest
    feasible inline expansion; patterns estimated to produce more
    than 50\,000 combinations are redirected to a conservative
    fallback via \texttt{\_process\_permute\_limited()}.
  \item \textbf{Permutation generation.}
    \texttt{ProductionPermuteHandler.expand\_permutation()} is called,
    which uses \texttt{itertools.per\-muta\-tions} for $n \leq 6$ and
    a memory-efficient streaming algorithm for larger $n$.  Results
    are memoised in a thread-safe LRU cache keyed on the sorted
    variable tuple.
  \item \textbf{NFA assembly.}  For each permutation
    $\sigma = (\sigma_1, \ldots, \sigma_n)$, the builder constructs
    a concatenation fragment by calling
    \texttt{create\_var\_states()} for each $\sigma_i$ and connecting
    the fragments with $\varepsilon$-transitions.  All $n!$
    concatenation fragments are then united under a common fork/merge
    pair with the same $\varepsilon$-alternation structure used for
    the \texttt{|} operator.
\end{enumerate}

% ................................
\subsection{Properties of the Constructed NFA}
\label{subsec:meth-nfa-properties}
% ................................

The NFA returned by \texttt{NFABuilder.build()} satisfies the
following formal properties:

\begin{enumerate}
  \item \textbf{Linear size.}  At most two states per
    \texttt{LITERAL} token (the transition's source and target)
    plus a bounded constant of intermediate states per structural
    token (alternation fork/merge, quantifier envelope, etc.).
    Total state count is $O(n)$ in token-list length $n$.

  \item \textbf{Single accept state.}  Exactly one accept state
    $q_\text{accept}$ exists; all fragment tails are connected to
    it via $\varepsilon$-transitions.

  \item \textbf{Priority-annotated $\varepsilon$-transitions.}
    Every $\varepsilon$-transition carries an integer priority
    stored in the \texttt{epsilon\_priorities} dictionary on its
    source \texttt{NFAState}.  Lower numeric values denote higher
    priority.  For greedy quantifiers the loop-back
    $\varepsilon$ has lower priority than the exit $\varepsilon$,
    so the $\varepsilon$-closure computation (which sorts targets by
    ascending priority) extends the match before terminating.
    Reluctant quantifiers reverse the ordering.  These priorities
    are propagated through subset construction so that each DFA
    transition inherits the minimum (highest-priority) value of the
    contributing NFA transitions, implementing SQL:2016 greedy and
    leftmost-match semantics without modifying the DFA construction
    algorithm.

  \item \textbf{Thread safety.}  Each \texttt{NFAState} and the
    \texttt{NFA} object itself are guarded by a
    \texttt{threading.RLock()}.  Epsilon-closure results are cached
    in a bounded dictionary (capped at 1\,000 entries) accessed
    under the same lock, enabling safe concurrent queries by
    multiple partition workers.

  \item \textbf{Structural validation.}  \texttt{NFA.validate()}
    verifies that (i)~all transition targets are valid state
    indices, (ii)~the accept state is reachable from the start
    state, and (iii)~PERMUTE metadata is internally consistent
    before the NFA is returned to the caller.
\end{enumerate}

% ................................
\subsection{Worked Example: Pattern \texttt{A\;B+\;C?}}
\label{subsec:meth-nfa-example}
% ................................

Consider the pattern \texttt{A B+ C?} with three pattern variables
\texttt{A}, \texttt{B}, \texttt{C} each associated with an
arbitrary Boolean DEFINE predicate.  The tokeniser produces the
flat list:
\begin{center}
  \texttt{[\,LITERAL('A'),\quad LITERAL('B', quantifier='+'),\quad
    LITERAL('C', quantifier='?')\,]}
\end{center}
The \texttt{NFABuilder} then processes this list as follows.

\begin{enumerate}
  \item A global start state $q_0$ and accept state $q_a$ are
    created; $q_a$ is marked \texttt{is\_accept = True}.

  \item \textbf{LITERAL \texttt{A} (no quantifier).}
    \texttt{create\_var\_states} allocates $q_1$ and $q_2$; adds
    transition $q_1 \xrightarrow{C(A)} q_2$.
    The visitor adds $\varepsilon(q_0 \to q_1)$ and advances
    \texttt{current} to $q_2$.

  \item \textbf{LITERAL \texttt{B} with quantifier \texttt{+}.}
    \texttt{create\_var\_states} allocates $q_3, q_4$ with
    $q_3 \xrightarrow{C(B)} q_4$.
    \texttt{\_apply\_quantifier}$(q_3,\,q_4,\,1,\,\infty,\,\mathit{True})$
    allocates envelope states $q_5, q_6$ and adds:
    $\varepsilon(q_5 \to q_3)$,\;
    $\varepsilon(q_4 \to q_6)$ (exit),\;
    $\varepsilon(q_4 \to q_3)$ (loop, greedy).
    The visitor adds $\varepsilon(q_2 \to q_5)$ and advances to $q_6$.

  \item \textbf{LITERAL \texttt{C} with quantifier \texttt{?}.}
    \texttt{create\_var\_states} allocates $q_7, q_8$ with
    $q_7 \xrightarrow{C(C)} q_8$.
    \texttt{\_apply\_quantifier}$(q_7,\,q_8,\,0,\,1,\,\mathit{True})$
    allocates $q_9, q_{10}$ and adds:
    $\varepsilon(q_9 \to q_7)$ (take $C$),\;
    $\varepsilon(q_8 \to q_{10})$ (exit after $C$),\;
    $\varepsilon(q_9 \to q_{10})$ (skip $C$, bypass).
    The visitor adds $\varepsilon(q_6 \to q_9)$ and advances to $q_{10}$.

  \item $\varepsilon(q_{10} \to q_a)$ connects the pattern's tail
    to the global accept state.
\end{enumerate}

\noindent
The resulting NFA (11 states: $q_0$–$q_{10}$ plus the accept state
shared with $q_{10}$ in this abbreviated numbering) accepts exactly
the language $A \cdot B^{+} \cdot C^{?}$: one row satisfying $A$,
one or more rows satisfying $B$ (the loop $q_4 \to q_3$ implements
the \texttt{+} greedily), and zero or one row satisfying $C$ (the
bypass $q_9 \to q_{10}$ implements the \texttt{?}).

After construction the NFA is passed to \texttt{DFABuilder.build()}
for subset construction as described in the next section.

% ................................
\subsection{Deterministic Finite Automata (DFA)}
\label{subsec:meth-dfa}
% ................................

A \emph{Deterministic Finite Automaton} (DFA) is a special case of
an NFA in which:
\begin{inparaenum}[(i)]
  \item there are no $\varepsilon$-transitions, and
  \item for each state $q$ and symbol $a$, $|\delta(q,a)| = 1$
    (exactly one successor state).
\end{inparaenum}
DFAs require $O(n)$ time to process an input of length $n$ with no
backtracking, making them ideal for row-by-row matching.

\paragraph{Subset Construction (NFA $\to$ DFA).}
The \emph{subset construction} (powerset construction) algorithm
converts an NFA to an equivalent DFA.  Each DFA state corresponds
to a \emph{set} of NFA states.  For a DFA state $S$ and input
symbol $a$:

\[
  \delta_{\text{DFA}}(S,\,a)
  \;=\;
  \varepsilon\text{-closure}\!\left(
    \bigcup_{q \in S} \delta_{\text{NFA}}(q,\,a)
  \right)
\]

where $\varepsilon\text{-closure}(T)$ is the set of all NFA states
reachable from any state in $T$ by zero or more
$\varepsilon$-transitions.  Although the resulting DFA may have up
to $2^{|Q|}$ states in the worst case, patterns arising from
\texttt{MATCH\_RECOGNIZE} syntax are structurally constrained and in
practice yield DFAs with far fewer states.

\paragraph{DFA Minimisation.}
Hopcroft's algorithm~\cite{HopcroftUllman} reduces a DFA to its
canonical minimal form in $O(n \log n)$ time by iteratively
partitioning states into equivalence classes of
indistinguishable states.  The implementation applies minimisation
after subset construction to reduce memory footprint and improve
matching speed.

\paragraph{State Prioritisation for SQL Semantics.}
SQL:2016 mandates \emph{greedy} matching by default: the engine
must prefer the longest possible match, and in the case of
alternation, the leftmost alternative takes precedence.  Standard
DFA theory does not encode these preferences.  The implementation
addresses this by:
\begin{enumerate}
  \item Assigning a \emph{priority} to each NFA transition based on
    its position in the pattern and quantifier greediness.
  \item Propagating priorities through subset construction so that
    each DFA transition inherits the minimum (highest-priority)
    value of the contributing NFA transitions.
  \item Ordering outgoing DFA transitions by priority during
    match execution so that higher-priority paths are explored
    first.
\end{enumerate}


% ----------------------------------------------------------------
\section{Pandas DataFrame Execution Model}
\label{sec:meth-pandas}
% ----------------------------------------------------------------

Pandas~\cite{McKinney2010} is a Python library providing the
\texttt{DataFrame} data structure: a two-dimensional, labelled
table backed by typed NumPy arrays.  Its in-memory, columnar
storage enables vectorised operations over millions of rows, making
it well suited for the row-by-row DFA simulation that
\texttt{MATCH\_RECOGNIZE} requires.

This section describes how the engine transforms a validated
\texttt{MATCH\_RECOGNIZE} clause and a compiled DFA into concrete
output rows.  The process passes through six distinct stages:
partition, order, row classification, DFA simulation,
measure evaluation, and output assembly.
Figure~\ref{fig:exec-pipeline} places these stages in context
of the full five-phase architecture introduced at the start of
this chapter.

\begin{figure}[htbp]
  \centering
  \begin{tikzpicture}[
    node distance=0.6cm and 1.6cm,
    every node/.style={font=\small},
    box/.style={rectangle, draw, rounded corners=3pt,
                minimum width=2.8cm, minimum height=0.75cm,
                align=center},
    arrow/.style={->, thick}
  ]
    \node[box, fill=gray!15] (parse) {SQL Parsing};
    \node[box, fill=gray!15, right=of parse] (sem) {Semantic\\Validation};
    \node[box, fill=gray!15, right=of sem] (plan) {Logical Plan};
    \node[box, fill=gray!15, right=of plan] (comp) {Pattern\\Compilation};
    \node[box, fill=blue!15, right=of comp] (exec) {\textbf{Execution}};
    \draw[arrow] (parse) -- (sem);
    \draw[arrow] (sem) -- (plan);
    \draw[arrow] (plan) -- (comp);
    \draw[arrow] (comp) -- (exec);
  \end{tikzpicture}
  \caption{Five-phase execution pipeline; the execution stage (shaded)
    is the subject of this section}
  \label{fig:exec-pipeline}
\end{figure}

% ................................
\subsection{Partition and Order}
\label{subsec:meth-exec-partition}
% ................................

Execution begins by splitting the input DataFrame $D$ into
independent \emph{partitions}.  The columns named in the
\texttt{PARTITION BY} clause are passed to
\texttt{DataFrame.groupby()}, which returns one sub-frame per
distinct combination of partition-key values.  When no
\texttt{PARTITION BY} clause is present the entire DataFrame
constitutes a single partition.

Each partition is then sorted by the \texttt{ORDER BY} columns
using \texttt{sort\_values()}, establishing the temporal order
within which the DFA simulation will proceed.  The underlying
NumPy index is preserved so that row references in navigation
functions remain stable.

Because partitions share no state, they can be processed
independently—and optionally in parallel (see
Section~\ref{subsec:bg-parallel}).  The remainder of this section
describes the processing of a single partition $P$ of $n$ rows,
indexed $r_0, r_1, \ldots, r_{n-1}$ in sorted order.

% ................................
\subsection{Row Classification}
\label{subsec:meth-exec-classify}
% ................................

Before the DFA simulation begins, every \texttt{DEFINE} predicate
is compiled into a callable Python function.  The
\texttt{compile\_condition()} function in
\texttt{src/matcher/condition\_evaluator.py} parses the predicate
expression—which may reference pattern variables via dot notation
(\texttt{A.price}), navigation functions, and SQL comparison
operators—and wraps it in a closure that accepts a row dictionary
and a \texttt{RowContext} object.

The \texttt{RowContext}
(Listing~\ref{lst:rowcontext-iface}) is the central evaluation
context.  It provides:
\begin{itemize}
  \item \texttt{rows}: the full sorted partition as a list of row
    dictionaries, enabling look-ahead and look-behind via index
    arithmetic;
  \item \texttt{variables}: the current match's variable bindings
    (a mapping from variable name to a list of matching row
    indices);
  \item \texttt{current\_idx}: the index of the row being tested;
  \item \texttt{current\_var}: the name of the pattern variable
    whose predicate is being evaluated; and
  \item \texttt{subsets}: the \texttt{SUBSET} variable definitions.
\end{itemize}

\begin{lstlisting}[
  language=Python,
  caption={Key fields of \texttt{RowContext} consulted during
    condition evaluation},
  label={lst:rowcontext-iface}]
@dataclass
class RowContext:
    rows:              List[Dict[str, Any]]   # full partition
    variables:         Dict[str, List[int]]   # var -> [row indices]
    current_idx:       int                    # row under test
    current_var:       Optional[str]          # variable being matched
    subsets:           Dict[str, List[str]]   # SUBSET definitions
    defined_variables: Set[str]               # DEFINE variable names
    pattern_variables: List[str]              # all pattern variables
\end{lstlisting}

\noindent
Predicates are evaluated lazily during DFA simulation: a condition
is called only when the DFA attempts the corresponding transition.
For large partitions (over $50\,000$ rows) the engine switches to a
vectorised evaluation path that pre-computes a
\emph{condition matrix}—a Boolean array of shape
$|\text{variables}| \times n$—using NumPy broadcast operations,
reducing condition evaluations from $O(n \cdot |\text{variables}|)$
scalar calls to $O(|\text{variables}|)$ vector operations.

% ................................
\subsection{DFA Simulation}
\label{subsec:meth-exec-dfa-sim}
% ................................

The core of the execution stage is the DFA simulation loop, shown
as Algorithm~\ref{alg:dfa-sim}.  Starting from the DFA's initial
state $q_0$, the engine processes the sorted partition one row at a
time.  At each row $r_i$ it evaluates the \texttt{DEFINE} predicates
for every outgoing transition of the current state, collects the
transitions whose predicates evaluate to \texttt{True}, and selects
the highest-priority one (reflecting greedy or leftmost-match
semantics).  When the DFA reaches an accepting state, the engine
records the current variable bindings as a candidate match and
either terminates (for simple patterns) or continues to seek a
longer greedy match.

\begin{algorithm}[htbp]
\caption{DFA Simulation Loop for a Single Partition}
\label{alg:dfa-sim}
\DontPrintSemicolon
\SetAlgoLined
\SetKwProg{Proc}{Procedure}{:}{}
\SetKw{KwRet}{return}
\Proc{\upshape\textsc{SimulatePartition}$(P,\, \mathit{dfa},\,
      \mathit{skip\_mode},\, \mathit{config})$}{
  $\mathit{results} \leftarrow [\,]$\;
  $i \leftarrow 0$\;
  \While{$i < |P|$}{
    $q \leftarrow q_0$\;
    $\mathit{bindings} \leftarrow \{\}$\;
    $\mathit{best} \leftarrow \bot$\;
    $j \leftarrow i$\;
    \While{$j < |P|$}{
      $\mathit{fired} \leftarrow
        \{(v,q') \mid (q \xrightarrow{v} q') \in \mathit{dfa},\;
          C(v, P[j], \mathit{ctx}_{j,\mathit{bindings}})\}$\;
      \uIf{$\mathit{fired} = \varnothing$}{
        \textbf{break}\tcp*{no transition -- end of match attempt}
      }
      $(v^*,q^*) \leftarrow \textsc{BestTransition}(\mathit{fired})$\;
      $\mathit{bindings}[v^*] \leftarrow
          \mathit{bindings}[v^*] \cup \{j\}$\;
      $q \leftarrow q^*$\;
      \uIf{$\mathit{dfa.accepting}(q)$}{
        $\mathit{best} \leftarrow
          (\mathit{start}{=}i,\;\mathit{end}{=}j,\;
           \mathit{bindings}\text{ copy})$\;
        \lIf{$\mathit{skip\_mode} = \textsc{ToNextRow}$}{\textbf{break}}
      }
      $j \leftarrow j + 1$\;
    }
    \uIf{$\mathit{best} \neq \bot$}{
      $\mathit{results}.\textsc{Append}(\mathit{best})$\;
      $i \leftarrow \textsc{NextStart}(i,\;\mathit{best},\;
        \mathit{skip\_mode},\;\mathit{bindings})$\;
    }\Else{
      $i \leftarrow i + 1$\;
    }
  }
  \KwRet $\mathit{results}$\;
}
\end{algorithm}

\paragraph{BestTransition selection.}
When multiple transitions fire simultaneously (e.g.\ when a row
satisfies more than one pattern variable's predicate) the engine
applies the following priority ordering, consistent with
SQL:2016 greedy-leftmost semantics:
\begin{enumerate}
  \item Transitions leading to accepting states are preferred over
    non-accepting ones, ensuring the shortest valid prefix is
    always recorded as a candidate match.
  \item Among non-accepting transitions, those whose source and
    target state are the same (self-loops for repeating quantifiers
    such as \texttt{A+}) are preferred over state-advancing
    transitions—implementing greedy extension.
  \item When the pattern contains cross-variable references in
    \texttt{DEFINE} predicates (e.g.\ \texttt{B.price > A.price}),
    the engine further biases toward the self-loop for the
    prerequisite variable (\texttt{A}) before switching to the
    dependent one (\texttt{B}), ensuring $A$'s rows are fully
    consumed before $B$ begins.
  \item Remaining ties are broken by alternation order: the
    leftmost alternative in the original \texttt{PATTERN} string
    receives the lowest numeric priority and is therefore selected
    first.
\end{enumerate}

\paragraph{Variable bindings.}
Throughout the simulation, \texttt{bindings} is a dictionary
$\{ v \mapsto [i_1, i_2, \ldots] \}$ mapping each pattern
variable $v$ to the (sorted) list of partition row indices
assigned to it during the current match attempt.  At the
conclusion of a successful match these lists provide the concrete
row sets over which measure expressions and navigation functions
are evaluated.

Rows belonging to \emph{exclusion} variables (those wrapped in
\texttt{\{- \ldots -\}}) are tracked in a separate
\texttt{excluded\_rows} list: they participate in condition
evaluation and DFA state transitions normally but are omitted from
the measure-evaluation context and from the output rows in
\texttt{ALL ROWS PER MATCH} mode, conforming to the SQL:2016
semantics for pattern exclusions.

% ................................
\subsection{After Match Skip}
\label{subsec:meth-exec-skip}
% ................................

After recording a match the engine must decide where to start the
next match attempt, controlling whether matches can overlap.
SQL:2016 defines four \texttt{AFTER MATCH SKIP} strategies;
Table~\ref{tab:skip-modes} summarises their semantics and the
corresponding next-start-position computation.

\begin{table}[htbp]
  \centering
  \caption{The four \texttt{AFTER MATCH SKIP} modes and the
    resulting next scan position}
  \label{tab:skip-modes}
  \begin{tabular}{lp{9cm}}
    \toprule
    \textbf{Mode} & \textbf{Next-start computation} \\
    \midrule
    \textsc{Past Last Row} (default) &
      $i \leftarrow \mathit{best.end} + 1$. No overlap is possible;
      the scan resumes immediately after the last matched row. \\[4pt]
    \textsc{To Next Row} &
      $i \leftarrow \mathit{best.start} + 1$. Overlapping matches
      are produced because only a single row is skipped regardless
      of match length. \\[4pt]
    \textsc{To First} $V$ &
      $i \leftarrow \min(\mathit{bindings}[V]) + 1$. The scan
      resumes after the \emph{first} row assigned to variable $V$.
      Requires $V$ to appear in the pattern; enforced at compile time. \\[4pt]
    \textsc{To Last} $V$ &
      $i \leftarrow \max(\mathit{bindings}[V]) + 1$. The scan
      resumes after the \emph{last} row assigned to variable $V$.
      May overlap if $V$'s last row precedes the match end. \\
    \bottomrule
  \end{tabular}
\end{table}

\noindent
To prevent infinite loops, the engine enforces a monotonic
progress invariant: the next scan position is always strictly
greater than the current start position.  For
\textsc{To Next Row} the engine additionally detects cycles in
the sequence of recent start positions and breaks out if the same
position is revisited.

% ................................
\subsection{Output Modes}
\label{subsec:meth-exec-output}
% ................................

SQL:2016 specifies two principal output-row multiplicity modes,
each admitting further sub-options
(Table~\ref{tab:rows-per-match}).

\begin{table}[htbp]
  \centering
  \caption{Output modes and their row-count and measure-semantics
    behaviour}
  \label{tab:rows-per-match}
  \begin{tabular}{llp{7cm}}
    \toprule
    \textbf{Mode} & \textbf{Rows out} & \textbf{Measure evaluation context} \\
    \midrule
    \textsc{One Row Per Match}
      & 1 per match
      & \textsc{Final}: measures see the complete, finalised
        variable bindings.  Navigation functions
        (\texttt{FIRST}, \texttt{LAST}) operate over the full
        matched span. \\[6pt]
    \textsc{All Rows Per Match}
      & 1 per matched row
      & Default semantics per function: \texttt{LAST} uses
        \textsc{Running}; \texttt{FIRST}/\texttt{PREV}/\texttt{NEXT}
        use \textsc{Final}; aggregates
        (\texttt{COUNT}, \texttt{SUM}) use \textsc{Running}.
        Explicit \texttt{RUNNING}/\texttt{FINAL} keywords override
        the default. \\[6pt]
    \textsc{All Rows \ldots With Unmatched}
      & 1 per every input row
      & Matched rows carry measure values; unmatched rows emit
        \texttt{NULL} for all measure columns. \\
    \bottomrule
  \end{tabular}
\end{table}

\paragraph{One Row Per Match.}
For each recorded match the engine constructs a single output
dictionary.  The \texttt{PARTITION BY} and \texttt{ORDER BY}
column values are copied from the first matched row; measure
expressions are evaluated once with \textsc{Final} semantics
(all variable bindings are visible).  The resulting dictionary is
appended to the results list and later assembled into the output
DataFrame.

\paragraph{All Rows Per Match.}
For each recorded match the engine iterates over every row
belonging to the match span (from \texttt{best.start} to
\texttt{best.end} inclusive).  For each row $r_j$ in the span
the engine:
\begin{enumerate}
  \item assigns $r_j$ a \texttt{CLASSIFIER} column value equal to
    the name of the pattern variable to which row $j$ was assigned
    (or \texttt{NULL} if the row was excluded);
  \item evaluates each measure expression with the semantics
    determined by Table~\ref{tab:rows-per-match}: \textsc{Running}
    semantics limits variable bindings to rows $\leq j$,
    while \textsc{Final} semantics always uses the complete
    bindings; and
  \item appends the resulting row dictionary to the results list.
\end{enumerate}
Rows covered by the match span but assigned to an exclusion
variable are handled according to the SQL:2016 rule: they appear
in the output with a \texttt{NULL} \texttt{CLASSIFIER} value and
\texttt{NULL} measure columns.

% ................................
\subsection{Measure Evaluation}
\label{subsec:meth-exec-measures}
% ................................

Measure expressions are evaluated by the
\texttt{MeasureEvaluator} class in
\texttt{src/matcher/measure\_evaluator.py}.  Each declared
measure is compiled into a callable similar to the condition
callables described in Section~\ref{subsec:meth-exec-classify}.
The evaluator dispatches on the expression type:

\begin{description}
  \item[Navigation functions.]
    \texttt{FIRST($v$.$c$)}, \texttt{LAST($v$.$c$)},
    \texttt{PREV($c$, $k$)}, and \texttt{NEXT($c$, $k$)}
    resolve to direct index lookups into the partition row list:
    \begin{itemize}
      \item $\texttt{FIRST}(v.c) \to
        P[\mathit{bindings}[v][0]][c]$
      \item $\texttt{LAST}(v.c) \to
        P[\mathit{bindings}[v][-1]][c]$
      \item $\texttt{PREV}(c,\,k) \to
        P[j - k][c]$ (where $j$ is the current output row index)
      \item $\texttt{NEXT}(c,\,k) \to
        P[j + k][c]$
    \end{itemize}
    All accesses are bounds-checked: out-of-partition references
    return \texttt{NULL} per SQL:2016.

  \item[Aggregate functions.]
    \texttt{COUNT}, \texttt{SUM}, and \texttt{AVG} are implemented
    as running or final aggregates over the row index sets stored
    in \texttt{bindings}.  Under \textsc{Running} semantics the
    aggregate is computed over
    $\{ i \in \mathit{bindings}[v] \mid i \leq j \}$;
    under \textsc{Final} semantics it uses all rows in
    $\mathit{bindings}[v]$.

  \item[Pattern-variable column references.]
    An expression of the form $v.c$ (without a function wrapper)
    resolves to the column $c$ of the \emph{first} row assigned to
    $v$.  In \textsc{Running} mode the reference is restricted to
    assigned rows up to the current output position.

  \item[Arithmetic and logical expressions.]
    Compound expressions (e.g.\ \texttt{LAST(A.price) -
    FIRST(A.price)}) are evaluated by recursively evaluating
    sub-expressions and applying the indicated operator.
\end{description}

\paragraph{RUNNING vs.\ FINAL semantics.}
The distinction between the two semantics is central to
\texttt{ALL ROWS PER MATCH} output.  Under \textsc{Running}
semantics a measure evaluated for output row $j$ sees only rows
$r_0, \ldots, r_j$—the match is \emph{in progress} and future
rows are unknown.  Under \textsc{Final} semantics the full
match is always available regardless of which row is being output.

Consider a \texttt{SUM(A.price)} measure over a match covering
rows 3, 5, and 7.  In \textsc{Running} mode the output for each
row is a cumulative partial sum:
\begin{center}
\begin{tabular}{cl}
  Output row & \textsc{Running} \texttt{SUM} \\
  \hline
  3 & $P[3].\mathit{price}$ \\
  5 & $P[3].\mathit{price} + P[5].\mathit{price}$ \\
  7 & $P[3].\mathit{price} + P[5].\mathit{price} + P[7].\mathit{price}$ \\
\end{tabular}
\end{center}
whereas \textsc{Final} mode emits the same total on all three
output rows.

% ................................
\subsection{Worked End-to-End Example}
\label{subsec:meth-exec-example}
% ................................

To illustrate the full execution pipeline consider the following
query applied to a stock-price DataFrame with columns
\texttt{symbol}, \texttt{date}, and \texttt{price}:

\begin{lstlisting}[language=SQL, caption={Example query: detecting
  a V-shape price pattern}, label={lst:exec-example-query}]
SELECT *
FROM prices
MATCH_RECOGNIZE (
  PARTITION BY symbol
  ORDER BY date
  MEASURES
    FIRST(A.price) AS start_price,
    LAST(B.price)  AS bottom_price,
    LAST(C.price)  AS end_price
  ONE ROW PER MATCH
  AFTER MATCH SKIP PAST LAST ROW
  PATTERN (A+ B+ C+)
  DEFINE
    A AS A.price >= PREV(A.price, 1),
    B AS B.price <  PREV(B.price, 1),
    C AS C.price >  PREV(C.price, 1)
)
\end{lstlisting}

\noindent
Suppose the partition for \texttt{symbol = 'AAPL'} contains
five rows sorted by \texttt{date} with prices
$100, 102, 98, 95, 103$.  Execution proceeds as follows.

\vspace{0.5em}
\noindent
\textbf{Step 1 — Partition and Order.}
\texttt{groupby('symbol')} produces a single-symbol sub-frame;
\texttt{sort\_values('date')} confirms the price sequence
$[100, 102, 98, 95, 103]$.

\vspace{0.5em}
\noindent
\textbf{Step 2 — DFA construction.}
The pattern \texttt{A+ B+ C+} compiles to an NFA with 8 states via
Thompson's construction and then to a DFA with 4 states via subset
construction (Section~\ref{subsec:meth-dfa}).

\vspace{0.5em}
\noindent
\textbf{Step 3 — Simulation from $i=0$.}

\begin{itemize}[leftmargin=2em]
  \item Row 0 ($p=100$): \texttt{PREV(A.price)} is undefined at the
    first row; the engine treats this as trivially satisfying
    $A$'s predicate.  Transition $q_0 \xrightarrow{A} q_1$ fires;
    $\mathit{bindings}[A] = [0]$.  State $q_1$ is not accepting.
  \item Row 1 ($p=102 \geq 100$): \texttt{A} predicate satisfied;
    self-loop $q_1 \xrightarrow{A} q_1$.  $\mathit{bindings}[A] = [0,1]$.
  \item Row 2 ($p=98 < 102$): \texttt{A} fails, \texttt{B} passes
    ($p < \mathit{prev}$).  Transition $q_1 \xrightarrow{B} q_2$.
    $\mathit{bindings}[B] = [2]$.
  \item Row 3 ($p=95 < 98$): \texttt{B} satisfied; self-loop
    $q_2 \xrightarrow{B} q_2$.  $\mathit{bindings}[B] = [2,3]$.
  \item Row 4 ($p=103 > 95$): \texttt{C} satisfied; transition
    $q_2 \xrightarrow{C} q_3$.  $q_3$ is accepting.
    $\mathit{best} = (\mathit{start}{=}0,\; \mathit{end}{=}4,\;
    \{A{:}[0,1],\; B{:}[2,3],\; C{:}[4]\})$.
\end{itemize}

\vspace{0.5em}
\noindent
\textbf{Step 4 — After Match Skip.}
Mode is \textsc{Past Last Row}: next scan starts at $i = 5$, which
equals $|P|$, so the simulation terminates.

\vspace{0.5em}
\noindent
\textbf{Step 5 — Measure evaluation (ONE ROW PER MATCH, FINAL).}
\begin{align*}
  \texttt{start\_price}  &= \texttt{FIRST}(A.\mathit{price})
    = P[0][\mathit{price}] = 100 \\
  \texttt{bottom\_price} &= \texttt{LAST}(B.\mathit{price})
    = P[3][\mathit{price}] = 95 \\
  \texttt{end\_price}    &= \texttt{LAST}(C.\mathit{price})
    = P[4][\mathit{price}] = 103
\end{align*}

\vspace{0.5em}
\noindent
\textbf{Step 6 — Output assembly.}
A single output row is emitted:
\texttt{(symbol='AAPL', start\_price=100, bottom\_price=95,
end\_price=103)}.

This example demonstrates how all execution components--partition
ordering, DFA simulation with greedy quantifiers, variable binding
accumulation, after-match skip, and measure evaluation—interact to
produce a single SQL:2016-compliant output row from a five-row
price sequence.


% ----------------------------------------------------------------
\section{Implementation and Evaluation Metrics}
\label{sec:meth-impl}
% ----------------------------------------------------------------

We implemented a \texttt{MATCH\_RECOGNIZE} engine for Pandas
DataFrames following the SQL:2016 standard for pattern matching in
sequential data.  The implementation has four core components: an
ANTLR4-based SQL parser that transforms \texttt{MATCH\_RECOGNIZE}
queries into abstract syntax trees; a finite-automata engine that
converts pattern expressions into optimised NFAs and DFAs using
Thompson's construction; a pattern-matching executor that processes
DataFrames with support for partitioning and ordering; and a measure
evaluator that handles navigation functions (\texttt{PREV},
\texttt{NEXT}, \texttt{FIRST}, \texttt{LAST}) and aggregations with
both running and final semantics.  The system supports advanced
pattern features including anchors, quantifiers, alternation,
permuted patterns, and exclusion syntax.

% ................................
\subsection{ANTLR4 Grammar and AST Construction}
\label{subsec:impl-antlr4}
% ................................

Our system uses a customised SQL grammar extended from Trino's SQL
that fully supports the SQL:2016 \texttt{MATCH\_RECOGNIZE} clause.
This grammar was defined using ANTLR4, covering complex patterns,
nested patterns, alternation, quantifiers, grouping, and
exclusions~\cite{antlr_trino_sql}.  The parser's output is a
detailed parse tree that captures the syntactic structure of each
SQL query~\cite{parr2013definitive,tarazona2021antlr}.

% ................................
\subsection{Finite Automata}
\label{subsec:impl-automata}
% ................................

\paragraph{Pattern Tokenisation.}
Pattern expressions are tokenised using a recursive-descent
tokenisation algorithm that supports anchors, grouping,
quantifiers, and nested constructs.

\paragraph{NFA Construction.}
We use Thompson's construction algorithm to translate pattern
tokens into a Non-deterministic Finite Automaton (NFA).  Each
state captures transitions, variable bindings, anchors, and
metadata for advanced features.  Quantifiers (\texttt{*},
\texttt{+}, \texttt{?}, \texttt{\{$n$,$m$\}}) are directly
mapped to sub-automata~\cite{georgiou2025streaming}.

\paragraph{DFA Construction and Optimisation.}
The NFA is systematically converted into a Deterministic Finite
Automaton (DFA) using subset construction, drastically reducing
the number of active states during execution.  Optimisation steps
including state minimisation, transition indexing, and merging of
equivalent states yield a compact, performant automaton.  The
resulting DFA supports efficient row-by-row evaluation while
preserving pattern semantics, priorities, and custom conditions
for accurate pattern
semantics~\cite{e104-d_3_381,10.1145/986537.986588,georgiou2025streaming}.

% ................................
\subsection{Execution}
\label{subsec:impl-execution}
% ................................

The matching engine applies the DFA sequentially over each partition
of the DataFrame with specialised handling for anchors, skip modes,
and edge cases.  The executor coordinates the entire
pattern-matching process, integrating parsing, automata
construction, match identification, and result formatting.

\paragraph{Condition and Measure Evaluation.}
A dedicated evaluation subsystem translates SQL pattern conditions
and measures into executable Python code that supports navigation
functions and aggregations with SQL:2016-compliant \texttt{NULL}
semantics for both running and final calculations.

% ----------------------------------------------------------------
\section{Performance Optimization Techniques}
\label{sec:bg-performance}
% ----------------------------------------------------------------

% ................................
\subsection{LRU Pattern Caching}
\label{subsec:bg-caching}
% ................................

Pattern compilation---tokenisation $\to$ NFA construction $\to$ DFA
conversion $\to$ minimisation---is computationally expensive,
especially for patterns with quantifiers and \texttt{PERMUTE}
constructs.  The implementation maintains a thread-safe
\emph{Least Recently Used} (LRU) cache that maps a hash of the
pattern string and \texttt{DEFINE} predicates to the compiled
pattern (the full NFA and DFA pipeline result).
Key configuration parameters are listed in
Table~\ref{tab:bg-cache-config}.

\begin{table}[htbp]
  \centering
  \caption{Pattern cache configuration defaults
    (\texttt{src/config/production\_config.py})}
  \label{tab:bg-cache-config}
  \begin{tabular}{lrl}
    \toprule
    \textbf{Parameter}         & \textbf{Default} & \textbf{Description} \\
    \midrule
    \texttt{cache\_size\_limit}        & 50\,000 entries  & Maximum cached patterns \\
    \texttt{cache\_memory\_limit\_mb}  & 2\,048 MB        & Maximum cache memory \\
    \texttt{cache\_ttl\_seconds}       & 3\,600 s         & Entry time-to-live \\
    \texttt{cache\_clear\_threshold\_mb} & 1\,600 MB      & Proactive eviction threshold \\
    \bottomrule
  \end{tabular}
\end{table}

% ................................
\subsection{Memory Management}
\label{subsec:bg-memory}
% ................................

Large-scale matching over millions of rows requires careful memory
management.  The implementation processes one partition at a time,
keeping only the current partition's data in memory at any point.
Memory usage is monitored via \texttt{psutil}; garbage collection
is triggered proactively when consumption approaches
user-configured limits (\texttt{cache\_clear\_threshold\_mb}).
An object-pool layer (\texttt{src/utils/memory\_management.py})
reuses frequently allocated objects across successive match
operations, reducing allocator pressure on long-running queries.

% ................................
\subsection{Vectorised Condition Evaluation}
\label{subsec:bg-vectorised}
% ................................

For large partitions the engine avoids repeated Python function-call
overhead by pre-computing a \emph{condition matrix}: a dictionary
mapping each pattern variable $v$ to a NumPy Boolean array of
length $n$ (the partition size), where entry $i$ is \texttt{True}
if row $i$ satisfies the \texttt{DEFINE} predicate for $v$.

The pre-computation pass in
\texttt{EnhancedMatcher.\_vectorize\_condition\_evaluation()}
evaluates each variable's predicate function once per row and
stores the results as a compact \texttt{numpy.ndarray(dtype=bool)}.
When the Polars library is available it is used to parallelise this
pass; otherwise a plain Python loop with NumPy storage provides the
same result.  Implicit variables (present in the pattern but absent
from \texttt{DEFINE}) are assigned an all-\texttt{True} array,
matching any row unconditionally.

Once the condition matrix is built, the DFA simulation loop
replaces every per-row predicate call with a single array index
look-up $M[v][j]$, reducing the inner-loop cost from an interpreted
function invocation to a direct memory access.

For partitions below $5\,000$ rows an enhanced
condition-evaluation cache is used instead; for very large
partitions ($>50\,000$ rows) a sample-based pre-warming step
warms internal caches before the full evaluation pass begins.
Both thresholds are hardcoded in
\texttt{src/matcher/matcher.py}.

Table~\ref{tab:perf-summary} summarises all performance
optimisations and the conditions under which each is applied.

\begin{table}[htbp]
  \centering
  \caption{Performance optimisations and their activation
    conditions}
  \label{tab:perf-summary}
  \begin{tabular}{p{3.8cm}p{4.2cm}p{4.4cm}}
    \toprule
    \textbf{Technique} & \textbf{Activation condition}
      & \textbf{Primary benefit} \\
    \midrule
    LRU pattern cache
      & Always on
      & Eliminates redundant NFA/DFA construction across repeated
        queries \\[4pt]
    Condition matrix (vectorised look-up)
      & Partition size $> 5\,000$ rows
      & Replaces per-row predicate calls with array index
        look-ups inside the DFA loop \\[4pt]
    One-partition-at-a-time streaming
      & Always on
      & Bounds peak memory to one partition regardless of total
        dataset size \\[4pt]
    Proactive GC / memory guard
      & Process RSS $\geq$ \texttt{cache\_clear\_threshold\_mb}
      & Prevents OOM errors on very large dataset scans \\
    \bottomrule
  \end{tabular}
\end{table}

% ----------------------------------------------------------------
\section{Chapter Summary}
\label{sec:meth-summary}
% ----------------------------------------------------------------

This chapter has presented the full design and implementation of a
SQL:2016-compliant \texttt{MATCH\_RECOGNIZE} engine built over
Pandas DataFrames.  The system is structured as a five-phase
pipeline:

\begin{enumerate}
  \item \textbf{SQL Parsing.}
    An ANTLR4-generated parser transforms a SQL query string into
    a structured AST.  The visitor pattern, implemented in
    \texttt{MatchRecognizeExtractor} and
    \texttt{FullQueryExtractor}, extracts all clause components
    while a safe-eval layer guards against code injection in
    user-supplied expressions.

  \item \textbf{Semantic Validation.}
    Before compilation begins, the engine verifies that all pattern
    variables referenced in \texttt{MEASURES} and \texttt{DEFINE}
    clauses are declared in \texttt{PATTERN}, that \texttt{SUBSET}
    definitions are consistent, and that navigation function
    arguments are type-compatible.

  \item \textbf{Logical Planning.}
    The validated AST is normalised into a canonical logical plan:
    implicit variable definitions (variables absent from
    \texttt{DEFINE}) are assigned an always-true predicate,
    skip-mode and output-mode defaults are applied, and
    \texttt{SUBSET} aliases are expanded.

  \item \textbf{Pattern Compilation.}
    The \texttt{PATTERN} string is tokenised into a typed token
    sequence; Thompson's construction converts it to an NFA; the
    subset-construction algorithm converts the NFA to a DFA; and
    Hopcroft's algorithm minimises the DFA.  The resulting DFA is
    stored in a thread-safe LRU cache keyed on a hash of the
    pattern and define predicates, ensuring that repeated
    executions of the same pattern incur no recompilation cost.

  \item \textbf{Execution.}
    The DataFrame is partitioned by \texttt{PARTITION BY} keys and
    each partition is sorted by \texttt{ORDER BY} columns.
    Partitions are executed in parallel when the partition count
    exceeds a configurable threshold.  Within each partition, the
    DFA simulation loop advances row by row, evaluating
    \texttt{DEFINE} predicates, accumulating variable bindings,
    recording matches at accepting states, and advancing the scan
    position according to the chosen \texttt{AFTER MATCH SKIP}
    strategy.  Measure expressions are evaluated under
    \textsc{Running} or \textsc{Final} semantics depending on the
    output mode and the function type.  The output is assembled
    into a new DataFrame with the declared measure columns.
\end{enumerate}

The design achieves correctness by adhering closely to the
SQL:2016 standard for pattern variable binding, greedy quantifier
semantics, after-match skip behaviour, and running/final measure
evaluation.  Performance at scale is addressed through LRU
compilation caching, parallel partition processing, and vectorised
condition evaluation, all of which are orthogonal to the core
matching logic and can be disabled individually for testing or
reproducibility purposes.

Chapter~5 evaluates these design decisions empirically, measuring
throughput, latency, and correctness against both synthetic
benchmarks and real-world financial event datasets.
